% Options for packages loaded elsewhere
\PassOptionsToPackage{unicode}{hyperref}
\PassOptionsToPackage{hyphens}{url}
%
\documentclass[
]{article}
\usepackage{amsmath,amssymb}
\usepackage{iftex}
\ifPDFTeX
  \usepackage[T1]{fontenc}
  \usepackage[utf8]{inputenc}
  \usepackage{textcomp} % provide euro and other symbols
\else % if luatex or xetex
  \usepackage{unicode-math} % this also loads fontspec
  \defaultfontfeatures{Scale=MatchLowercase}
  \defaultfontfeatures[\rmfamily]{Ligatures=TeX,Scale=1}
\fi
\usepackage{lmodern}
\ifPDFTeX\else
  % xetex/luatex font selection
\fi
% Use upquote if available, for straight quotes in verbatim environments
\IfFileExists{upquote.sty}{\usepackage{upquote}}{}
\IfFileExists{microtype.sty}{% use microtype if available
  \usepackage[]{microtype}
  \UseMicrotypeSet[protrusion]{basicmath} % disable protrusion for tt fonts
}{}
\makeatletter
\@ifundefined{KOMAClassName}{% if non-KOMA class
  \IfFileExists{parskip.sty}{%
    \usepackage{parskip}
  }{% else
    \setlength{\parindent}{0pt}
    \setlength{\parskip}{6pt plus 2pt minus 1pt}}
}{% if KOMA class
  \KOMAoptions{parskip=half}}
\makeatother
\usepackage{xcolor}
\usepackage{color}
\usepackage{fancyvrb}
\newcommand{\VerbBar}{|}
\newcommand{\VERB}{\Verb[commandchars=\\\{\}]}
\DefineVerbatimEnvironment{Highlighting}{Verbatim}{commandchars=\\\{\}}
% Add ',fontsize=\small' for more characters per line
\newenvironment{Shaded}{}{}
\newcommand{\AlertTok}[1]{\textcolor[rgb]{1.00,0.00,0.00}{\textbf{#1}}}
\newcommand{\AnnotationTok}[1]{\textcolor[rgb]{0.38,0.63,0.69}{\textbf{\textit{#1}}}}
\newcommand{\AttributeTok}[1]{\textcolor[rgb]{0.49,0.56,0.16}{#1}}
\newcommand{\BaseNTok}[1]{\textcolor[rgb]{0.25,0.63,0.44}{#1}}
\newcommand{\BuiltInTok}[1]{\textcolor[rgb]{0.00,0.50,0.00}{#1}}
\newcommand{\CharTok}[1]{\textcolor[rgb]{0.25,0.44,0.63}{#1}}
\newcommand{\CommentTok}[1]{\textcolor[rgb]{0.38,0.63,0.69}{\textit{#1}}}
\newcommand{\CommentVarTok}[1]{\textcolor[rgb]{0.38,0.63,0.69}{\textbf{\textit{#1}}}}
\newcommand{\ConstantTok}[1]{\textcolor[rgb]{0.53,0.00,0.00}{#1}}
\newcommand{\ControlFlowTok}[1]{\textcolor[rgb]{0.00,0.44,0.13}{\textbf{#1}}}
\newcommand{\DataTypeTok}[1]{\textcolor[rgb]{0.56,0.13,0.00}{#1}}
\newcommand{\DecValTok}[1]{\textcolor[rgb]{0.25,0.63,0.44}{#1}}
\newcommand{\DocumentationTok}[1]{\textcolor[rgb]{0.73,0.13,0.13}{\textit{#1}}}
\newcommand{\ErrorTok}[1]{\textcolor[rgb]{1.00,0.00,0.00}{\textbf{#1}}}
\newcommand{\ExtensionTok}[1]{#1}
\newcommand{\FloatTok}[1]{\textcolor[rgb]{0.25,0.63,0.44}{#1}}
\newcommand{\FunctionTok}[1]{\textcolor[rgb]{0.02,0.16,0.49}{#1}}
\newcommand{\ImportTok}[1]{\textcolor[rgb]{0.00,0.50,0.00}{\textbf{#1}}}
\newcommand{\InformationTok}[1]{\textcolor[rgb]{0.38,0.63,0.69}{\textbf{\textit{#1}}}}
\newcommand{\KeywordTok}[1]{\textcolor[rgb]{0.00,0.44,0.13}{\textbf{#1}}}
\newcommand{\NormalTok}[1]{#1}
\newcommand{\OperatorTok}[1]{\textcolor[rgb]{0.40,0.40,0.40}{#1}}
\newcommand{\OtherTok}[1]{\textcolor[rgb]{0.00,0.44,0.13}{#1}}
\newcommand{\PreprocessorTok}[1]{\textcolor[rgb]{0.74,0.48,0.00}{#1}}
\newcommand{\RegionMarkerTok}[1]{#1}
\newcommand{\SpecialCharTok}[1]{\textcolor[rgb]{0.25,0.44,0.63}{#1}}
\newcommand{\SpecialStringTok}[1]{\textcolor[rgb]{0.73,0.40,0.53}{#1}}
\newcommand{\StringTok}[1]{\textcolor[rgb]{0.25,0.44,0.63}{#1}}
\newcommand{\VariableTok}[1]{\textcolor[rgb]{0.10,0.09,0.49}{#1}}
\newcommand{\VerbatimStringTok}[1]{\textcolor[rgb]{0.25,0.44,0.63}{#1}}
\newcommand{\WarningTok}[1]{\textcolor[rgb]{0.38,0.63,0.69}{\textbf{\textit{#1}}}}
\usepackage{longtable,booktabs,array}
\usepackage{calc} % for calculating minipage widths
% Correct order of tables after \paragraph or \subparagraph
\usepackage{etoolbox}
\makeatletter
\patchcmd\longtable{\par}{\if@noskipsec\mbox{}\fi\par}{}{}
\makeatother
% Allow footnotes in longtable head/foot
\IfFileExists{footnotehyper.sty}{\usepackage{footnotehyper}}{\usepackage{footnote}}
\makesavenoteenv{longtable}
\setlength{\emergencystretch}{3em} % prevent overfull lines
\providecommand{\tightlist}{%
  \setlength{\itemsep}{0pt}\setlength{\parskip}{0pt}}
\setcounter{secnumdepth}{-\maxdimen} % remove section numbering
\ifLuaTeX
  \usepackage{selnolig}  % disable illegal ligatures
\fi
\usepackage{bookmark}
\IfFileExists{xurl.sty}{\usepackage{xurl}}{} % add URL line breaks if available
\urlstyle{same}
\hypersetup{
  hidelinks,
  pdfcreator={LaTeX via pandoc}}

\author{}
\date{}

\begin{document}

\section{A Ledger-Gated Constructive RG Program for Three-Dimensional
Yang--Mills
Theory}\label{a-ledger-gated-constructive-rg-program-for-three-dimensional-yangmills-theory}

\textbf{Version:} 1.8\\
\textbf{Date:} 2026-05-28\\
\textbf{Author:} Benjamin John Schulz\\
\textbf{License:} MIT\\
\textbf{Note:} This PDF/\LaTeX{} rendering tracks the Markdown source. Version~1.8 adds an open-source orchestration harness and seven structural improvements (single authoritative gate graph, explicit kill/pivot budget, physically enforced checker independence, an automated claim ratchet, resource schedules, a pinned exact-arithmetic/determinism policy, and schema-discriminator repairs); these are documented in full in Section~20 of \texttt{YM3D\_CONSTRUCTIVE\_RG\_DESIGN.md} and shipped as code in the repository.

\subsection{Abstract}\label{abstract}

This document specifies a ledger-gated constructive
renormalization-group (RG) program for three-dimensional Euclidean
lattice Yang--Mills theory, initially for compact gauge group (SU(2))
with Wilson action. The program is not presented as a completed proof.
It is a formal architecture for converting local block-RG calculations
into auditable mathematical claims, with exact arithmetic, hard
acceptance gates, scale-indexed polymer activity bounds, block-gauge
metadata, small/large-field partition ledgers, and periodic
thermodynamic-limit compatibility checks. Version 1.7 is the
validation-lane and claim-anchoring revision: it preserves the v1.6
implementation-readiness gates and adds an exact two-dimensional SU(N)
validation lane, a numerical anchor lane, an observable registry focused
on gauge-invariant Wilson-loop and glueball observables, explicit
claim-strength annotations in the pipeline DAG, expanded
primitive-enclosure backends, and a staged formal finite-group audit
lane after Milestone 1A.

The design is motivated by a successful micro-lemma pattern extracted
from recent K-star terminal-factor work:

{[} \text{local atoms} \to \text{collision graph} \to
\text{absolute majorant} \to \text{cluster/Neumann control} \to
\text{periodic lift} \to \text{KP activity map} \to
\text{explicit residual budget}. {]}

The central thesis is that this pattern can be transferred to
three-dimensional Yang--Mills only after several 3D-specific corrections
are made: all activity bounds must be scale/coupling-indexed; all atoms
must carry gauge-lane and reflection-positivity metadata; all Taylor
terms must carry small/large-field cutoff data; local geometry must be
symmetry-reduced before full expansion; block-gauge scale-transition
ripples must enter the coupling-flow error budget; and large-field
defect suppression must strictly dominate the relevant vector-valued
defect entropy rate.

Version 1.2 adds an executable-spine discipline: no additional ledger
family is introduced unless it blocks an acceptance gate; terminal
K-star closure exports a typed activity vector rather than an RG rung;
polymer convergence criteria are modular; fallback-basin certification
is separated from UV descent; and comparison lanes for scalar, abelian,
stochastic, and Hamiltonian approaches are advisory rather than
dependencies.

Version 1.3 is a schema-freeze readiness revision. It repairs the
reflection-positivity transfer pipeline, fixes dyadic blocking as the
default production scale, restricts tree/block-axial gauge lanes to
local UV and intermediate-scale coordinate use, generalizes large-field
defects from scalar surface area to vector-valued support measures,
makes the lattice-spacing-uniform gap milestone explicitly dependent on
Gate S13, and specifies the cubic group implementation by signed
permutation matrices with an independent link-coordinate audit.

Version 1.4 adds primitive enclosure certificates, reflection-plane
factorization checks, determinant-expansion hooks, boundary-surface
multiplier certificates, and mixed-defect junction entropy metadata.
Version 1.5 removes the null-field schema risk by requiring
discriminated unions keyed by \texttt{schema\_stage} or
\texttt{milestone\_target}; it also defines composite support action
rules, determinant/ghost measure hooks, norm-typed flow errors, and the
default order of limits. Version 1.6 adds the last
implementation-readiness controls before Milestone 1A: orbit-closure
consistency, ghost determinant sign admissibility, a diagrammatic and
machine-readable pipeline DAG, and schema-discriminator null tests.
Version 1.7 adds validation lanes and claim-anchoring controls: exact
two-dimensional SU(N) validation before analytic rows, numerical anchor
checks, an observable registry, claim-strength DAG annotations, and an
earlier formal finite-group audit lane.

\begin{center}\rule{0.5\linewidth}{0.5pt}\end{center}

\subsection{1. Scope and Claim
Boundaries}\label{scope-and-claim-boundaries}

\subsubsection{1.1 Target model}\label{target-model}

The initial target is three-dimensional Euclidean lattice Yang--Mills
theory on a cubic lattice with compact gauge group (SU(2)). The baseline
measure is the Wilson lattice measure

{[} Z\_\Lambda(\beta) = \int \prod\emph{\{e\in E(\Lambda)\} dU\_e
\exp\left( \frac{\beta}{2}\sum}\{p\in P(\Lambda)\}
\operatorname{Re}\operatorname{Tr} U\_p \right), {]}

where (dU\_e) is normalized Haar measure on (SU(2)), (U\_p) is the
plaquette holonomy, and (\Lambda) is a finite periodic cubic box unless
otherwise stated.

The physical three-dimensional coupling (g\_3\^{}2) has mass dimension
one. In the Wilson normalization used here, the weak-coupling scaling
has the schematic form

{[} \beta(a)\asymp \frac{1}{g_3^2 a}. {]}

Thus, under dyadic coarse-graining (a\_\{k+1\}=2a\_k), the leading
dimensional flow is expected to move (\beta\_k) toward the
strong-coupling/fallback regime.

\subsubsection{1.2 Near-term theorem
target}\label{near-term-theorem-target}

The near-term target is not the full continuum theorem. The first
serious theorem target is a scale-indexed one-step RG rung:

{[} \mathcal N\_\{\rm KP\}\^{}\{(k+1)\}(S'\emph{\{\rm eff\}) \le
\Theta\emph{k,\mathcal N}\{\rm KP\}\^{}\{(k)\}(S}\{\rm eff\}) +E\_k,
\qquad \Theta\_k\textless1, {]}

or a finite-step controlled descent into a certified fallback basin.

\subsubsection{1.3 Longer-term theorem
target}\label{longer-term-theorem-target}

The longer-term target is a volume-uniform and lattice-spacing-uniform
clustering estimate for gauge-invariant local observables (O):

{[} \left\textbar{}\langle O(x)O(y)\rangle\emph{c\right\textbar{} \le
C\_O e\^{}\{-m}\{\rm phys\}\textbar x-y\textbar\}, \qquad
m\_\{\rm phys\}\ge c g\_3\^{}2, \qquad c\textgreater0, {]}

with constants that survive the continuum scaling window.

\subsubsection{1.4 Non-claims}\label{non-claims}

This document does not claim:

\begin{enumerate}
\def\labelenumi{\arabic{enumi}.}
\tightlist
\item
  a completed construction of continuum Yang--Mills theory;
\item
  a completed proof of a mass gap;
\item
  a completed Balaban-style RG proof;
\item
  validity of any uncomputed K-star or nonterminal remainder estimate;
\item
  use of signed cancellation unless separately certified;
\item
  reflection-positivity transfer for gauge-fixed or smeared observables
  unless separately certified.
\end{enumerate}

\subsubsection{1.5 Anti-bureaucracy rule and minimal executable
spine}\label{anti-bureaucracy-rule-and-minimal-executable-spine}

The ledger machinery is intentionally strict, but it must not become an
independent bureaucracy. A new ledger family may be added only if it
blocks an acceptance gate, closes a named milestone, or supplies an
independent negative control for an already-declared gate.

The first implementation target is the minimal viable ledger spine

{[} \mathrm{MVL}\text{-}1 = \mathrm{geometry} + \mathrm{gauge\ trees} +
\mathrm{orbit\ representatives} + \mathrm{expanded\ graph\ audit}. {]}

The associated row sequence is

\begin{Shaded}
\begin{Highlighting}[]
\NormalTok{YM3D\_SCHEMA\_DISCRIMINATOR\_NULL\_TEST\_1}
\NormalTok{YM3D\_LOCAL\_ATOM\_GEOMETRY\_1}
\NormalTok{YM3D\_BLOCK\_GAUGE\_FIXING\_LEDGER\_1}
\NormalTok{YM3D\_BALABAN\_TREE\_AXIAL\_GAUGE\_HOOK\_1}
\NormalTok{YM3D\_HAAR\_REVERSIBILITY\_GAUGE\_CHECK\_1}
\NormalTok{YM3D\_CUBIC\_SYMMETRY\_ORBIT\_REDUCTION\_1}
\NormalTok{YM3D\_ORBIT\_CLOSURE\_CONSISTENCY\_CHECK\_1}
\NormalTok{YM3D\_ORBIT\_REDUCTION\_EXPANDED\_GRAPH\_AUDIT\_1}
\end{Highlighting}
\end{Shaded}

No Taylor, BKAR, defect-tail, continuum-limit, or mass-gap claim is part
of MVL-1.

\subsubsection{1.6 Schema-stage discipline and discriminated
unions}\label{schema-stage-discipline-and-discriminated-unions}

Version 1.5 replaces nullable analytic fields with stage-discriminated
schema branches. A row must declare either a \texttt{schema\_stage} or a
\texttt{milestone\_target}, and the JSON Schema must use \texttt{oneOf}
branches keyed by that discriminator.

The baseline stages are:

\begin{Shaded}
\begin{Highlighting}[]
\NormalTok{combinatorial}
\NormalTok{gauge\_geometric}
\NormalTok{analytic}
\NormalTok{rp\_transfer}
\NormalTok{continuum\_uniformity}
\end{Highlighting}
\end{Shaded}

A combinatorial Milestone 1A row must not merely set analytic fields to
\texttt{null}; those fields are forbidden unless a later stage branch
explicitly requires them. Conversely, an analytic row that invokes a
primitive constant, determinant expansion, cutoff boundary, or RP
transfer must populate the corresponding certificate identifiers.

The design therefore separates responsibilities:

\begin{longtable}[]{@{}
  >{\raggedright\arraybackslash}p{(\columnwidth - 2\tabcolsep) * \real{0.5000}}
  >{\raggedright\arraybackslash}p{(\columnwidth - 2\tabcolsep) * \real{0.5000}}@{}}
\toprule\noalign{}
\begin{minipage}[b]{\linewidth}\raggedright
Layer
\end{minipage} & \begin{minipage}[b]{\linewidth}\raggedright
Responsibility
\end{minipage} \\
\midrule\noalign{}
\endhead
\bottomrule\noalign{}
\endlastfoot
JSON Schema & structural typing, required/forbidden fields, enums,
rational interval shape \\
Python generator & canonical row and artifact generation \\
Independent Python checker & mathematical validation: orbits,
stabilizers, graph counts, support intersections, interval arithmetic,
DAG dependencies \\
Later proof-assistant lane & formalization of stable finite
combinatorics after the object model stabilizes \\
\end{longtable}

A typical structural pattern is:

\begin{Shaded}
\begin{Highlighting}[]
\FunctionTok{\{}
  \DataTypeTok{"oneOf"}\FunctionTok{:} \OtherTok{[}
    \FunctionTok{\{}
      \DataTypeTok{"properties"}\FunctionTok{:} \FunctionTok{\{}\DataTypeTok{"schema\_stage"}\FunctionTok{:} \FunctionTok{\{}\DataTypeTok{"const"}\FunctionTok{:} \StringTok{"combinatorial"}\FunctionTok{\}\},}
      \DataTypeTok{"required"}\FunctionTok{:} \OtherTok{[}\StringTok{"schema\_stage"}\OtherTok{,} \StringTok{"support\_id"}\OtherTok{,} \StringTok{"anchor\_type"}\OtherTok{,} \StringTok{"orientation"}\OtherTok{,} \StringTok{"orbit\_rep\_id"}\OtherTok{]}\FunctionTok{,}
      \DataTypeTok{"not"}\FunctionTok{:} \FunctionTok{\{}
        \DataTypeTok{"anyOf"}\FunctionTok{:} \OtherTok{[}
          \FunctionTok{\{}\DataTypeTok{"required"}\FunctionTok{:} \OtherTok{[}\StringTok{"primitive\_enclosure\_ids"}\OtherTok{]}\FunctionTok{\}}\OtherTok{,}
          \FunctionTok{\{}\DataTypeTok{"required"}\FunctionTok{:} \OtherTok{[}\StringTok{"boundary\_surface\_multiplier\_id"}\OtherTok{]}\FunctionTok{\}}\OtherTok{,}
          \FunctionTok{\{}\DataTypeTok{"required"}\FunctionTok{:} \OtherTok{[}\StringTok{"determinant\_expansion\_id"}\OtherTok{]}\FunctionTok{\}}
        \OtherTok{]}
      \FunctionTok{\}}
    \FunctionTok{\}}\OtherTok{,}
    \FunctionTok{\{}
      \DataTypeTok{"properties"}\FunctionTok{:} \FunctionTok{\{}\DataTypeTok{"schema\_stage"}\FunctionTok{:} \FunctionTok{\{}\DataTypeTok{"const"}\FunctionTok{:} \StringTok{"analytic"}\FunctionTok{\}\},}
      \DataTypeTok{"required"}\FunctionTok{:} \OtherTok{[}\StringTok{"schema\_stage"}\OtherTok{,} \StringTok{"support\_id"}\OtherTok{,} \StringTok{"coefficient\_bound"}\OtherTok{,} \StringTok{"norm\_id"}\OtherTok{,} \StringTok{"primitive\_enclosure\_ids"}\OtherTok{]}
    \FunctionTok{\}}
  \OtherTok{]}
\FunctionTok{\}}
\end{Highlighting}
\end{Shaded}

\subsubsection{1.7 Schema-discriminator null
test}\label{schema-discriminator-null-test}

Before the Milestone 1A generator emits any geometric atom, the
repository must run a discriminator null-test harness:

\begin{Shaded}
\begin{Highlighting}[]
\NormalTok{YM3D\_SCHEMA\_DISCRIMINATOR\_NULL\_TEST\_1}
\NormalTok{validate\_schema\_discriminators.py}
\end{Highlighting}
\end{Shaded}

The null-test verifies that the \texttt{oneOf}/\texttt{anyOf}
discriminator branches are not merely documentary. At minimum it
contains four fixtures:

\begin{longtable}[]{@{}
  >{\raggedright\arraybackslash}p{(\columnwidth - 2\tabcolsep) * \real{0.5000}}
  >{\raggedright\arraybackslash}p{(\columnwidth - 2\tabcolsep) * \real{0.5000}}@{}}
\toprule\noalign{}
\begin{minipage}[b]{\linewidth}\raggedright
Fixture
\end{minipage} & \begin{minipage}[b]{\linewidth}\raggedright
Expected result
\end{minipage} \\
\midrule\noalign{}
\endhead
\bottomrule\noalign{}
\endlastfoot
valid combinatorial row with only combinatorial fields & pass \\
combinatorial row containing analytic-only fields such as
\texttt{primitive\_enclosure\_ids} & fail \\
analytic row missing required analytic certificate fields & fail \\
row with an invalid discriminator value such as
\texttt{schema\_stage\ =\ partial} & fail \\
\end{longtable}

A failed null-test blocks Milestone 1A generation. This is a structural
test, not a mathematical proof, but it prevents the generator from
building expensive data on top of permissive or ambiguous schemas.

\subsubsection{1.8 Pipeline DAG
discipline}\label{pipeline-dag-discipline}

The proof pipeline must be represented both in prose and as
machine-readable artifacts:

\begin{Shaded}
\begin{Highlighting}[]
\NormalTok{docs/pipeline\_dag.mmd}
\NormalTok{artifacts/pipeline\_dag.json}
\NormalTok{schemas/pipeline\_dag.schema.json}
\end{Highlighting}
\end{Shaded}

The conceptual pipeline is:

\begin{Shaded}
\begin{Highlighting}[]
\NormalTok{flowchart TD}
\NormalTok{    A[Local Atomic Geometry] {-}{-}\textgreater{} B[Gauge{-}Lane Assignment]}
\NormalTok{    B {-}{-}\textgreater{} C[O\_h Link Action]}
\NormalTok{    C {-}{-}\textgreater{} D[Upward Support Composition]}
\NormalTok{    D {-}{-}\textgreater{} E[Orbit Representatives]}
\NormalTok{    E {-}{-}\textgreater{} F[Orbit{-}Closure Consistency]}
\NormalTok{    F {-}{-}\textgreater{} G[Expanded{-}Graph Audit]}
\NormalTok{    G {-}{-}\textgreater{} H[Scale/Coupling Slabs]}
\NormalTok{    H {-}{-}\textgreater{} I[Terminal Activity Vector]}
\NormalTok{    I {-}{-}\textgreater{} J[Majorant Matrix]}
\NormalTok{    J {-}{-}\textgreater{} K[Polymer Convergence Gate]}

\NormalTok{    H {-}{-}\textgreater{} L[Small/Large{-}Field Partition]}
\NormalTok{    L {-}{-}\textgreater{} M[Taylor Bucket]}
\NormalTok{    L {-}{-}\textgreater{} N[Tail/Defect Bucket]}
\NormalTok{    M {-}{-}\textgreater{} O[Nonterminal Inventory]}
\NormalTok{    N {-}{-}\textgreater{} O}
\NormalTok{    O {-}{-}\textgreater{} K}

\NormalTok{    B {-}{-}\textgreater{} P[Haar/RP Transfer]}
\NormalTok{    P {-}{-}\textgreater{} K}

\NormalTok{    O {-}{-}\textgreater{} Q[Ghost/Determinant Sign Check]}
\NormalTok{    Q {-}{-}\textgreater{} K}
\end{Highlighting}
\end{Shaded}

The machine-readable DAG records node identifiers, artifact types,
required parent nodes, allowed schema stages, checker entry points, and
claims unlocked by successful validation. A row that is not reachable
from the declared DAG is advisory and may not be imported as a proof
dependency without a transfer row.

\subsubsection{1.9 Exact validation and numerical anchor
policy}\label{exact-validation-and-numerical-anchor-policy}

The schema/generator pipeline must be tested on solved or externally
constrained cases before analytic three-dimensional rows are allowed to
close.

The primary validation lane is

\begin{Shaded}
\begin{Highlighting}[]
\NormalTok{YM2D\_SUN\_EXACT\_VALIDATION\_LANE\_0}
\end{Highlighting}
\end{Shaded}

This lane is not a prerequisite for the purely combinatorial Milestone
1A generator, because Milestone 1A is three-dimensional geometry and
symmetry infrastructure. It is, however, a prerequisite for any analytic
row claiming terminal KP smallness, Taylor/tail remainder control,
clustering, or fallback-basin certification. The purpose is to run the
same ledger machinery on two-dimensional SU(N) Yang--Mills, where exact
Wilson-loop and area-law formulas are available for comparison.

A second advisory lane is

\begin{Shaded}
\begin{Highlighting}[]
\NormalTok{YM3D\_NUMERICAL\_ANCHOR\_LANE\_0}
\end{Highlighting}
\end{Shaded}

This lane compares generated terminal activity vectors, Taylor
coefficients, Wilson-loop estimates, and glueball-correlator proxies
against high-precision small-volume or public lattice numerical data.
Numerical anchors never close theorem gates by themselves. They produce
source-audit warnings when analytic coefficients, normalization
conventions, or gauge-ripple budgets disagree with independent numerical
evidence.

\subsubsection{1.10 Observable-driven claim
discipline}\label{observable-driven-claim-discipline}

Every clustering, fallback, or continuum-uniformity claim must specify
the gauge-invariant observables to which it applies. The project
therefore maintains

\begin{Shaded}
\begin{Highlighting}[]
\NormalTok{YM3D\_OBSERVABLE\_REGISTRY\_1}
\end{Highlighting}
\end{Shaded}

with initial observable classes:

\begin{longtable}[]{@{}
  >{\raggedright\arraybackslash}p{(\columnwidth - 2\tabcolsep) * \real{0.5000}}
  >{\raggedright\arraybackslash}p{(\columnwidth - 2\tabcolsep) * \real{0.5000}}@{}}
\toprule\noalign{}
\begin{minipage}[b]{\linewidth}\raggedright
Observable
\end{minipage} & \begin{minipage}[b]{\linewidth}\raggedright
Role
\end{minipage} \\
\midrule\noalign{}
\endhead
\bottomrule\noalign{}
\endlastfoot
Wilson loop \texttt{W(C)} & area/perimeter-law diagnostics and exact 2D
validation \\
scalar glueball operator \texttt{O\_\{0++\}} & mass-gap/clustering
target \\
connected scalar glueball correlator \texttt{G\_\{0++\}(x,y)} &
quantitative exponential-clustering target \\
Polyakov loop \texttt{P(x)} & finite-temperature diagnostic only, not a
zero-temperature mass-gap observable \\
\end{longtable}

Gate S13 must ultimately connect lattice-spacing-uniformity to an
observable-specific statement such as

{[} G\_\{0\^{}\{++\}\}(x,y) \le C
e\^{}\{-m\_\{\rm phys\}\textbar x-y\textbar\}, \qquad
m\_\{\rm phys\}\ge c g\_3\^{}2, \qquad c\textgreater0, {]}

with constants uniform over the declared scaling window. Until such an
observable-specific certificate exists, the project may claim only the
weaker lattice or volume-uniform bounds explicitly unlocked by earlier
gates.

\begin{center}\rule{0.5\linewidth}{0.5pt}\end{center}

\subsection{2. Lessons Extracted from the K-Star Micro-Lemma
Chain}\label{lessons-extracted-from-the-k-star-micro-lemma-chain}

\subsubsection{2.1 A formula is not a result until it closes a ledger
row}\label{a-formula-is-not-a-result-until-it-closes-a-ledger-row}

The micro-lemma construction became mathematically useful only when
local formulas were converted into rows with:

\begin{itemize}
\tightlist
\item
  exact scope;
\item
  parent-row dependencies;
\item
  exact arithmetic;
\item
  machine-checkable acceptance criteria;
\item
  negative controls;
\item
  explicit non-claims;
\item
  and a named next obligation.
\end{itemize}

For the 3D program, a row is considered closed only if it answers:

{[} \text{What precise theorem fragment is now proved?} {]}

and

{[} \text{What precise theorem fragment remains open?} {]}

\subsubsection{2.2 Terminal and nonterminal activity must be
separated}\label{terminal-and-nonterminal-activity-must-be-separated}

The K-star terminal factor closed because it was sharply scoped. It was
possible to enumerate terminal atoms, form a collision graph, majorize
it, lift it to periodic boxes, and place it inside a KP/Balaban activity
map. The nonterminal remainder did not close until a source inventory
existed.

Therefore the 3D program must maintain two lanes:

{[} \mathcal A\_\{\rm terminal\}\^{}\{3D\}(k,I\_k) {]}

and

{[} \mathcal R\_\{\rm nonterminal\}\^{}\{3D\}(k,I\_k). {]}

Terminal success must never be promoted to full-RG success.

\subsubsection{2.3 Sign balance is diagnostic, not
foundational}\label{sign-balance-is-diagnostic-not-foundational}

A sign-orientation ledger may reveal structure, but no signed
cancellation is usable unless a signed analytic covariance or signed
mixed-term estimate is separately imported and checked. The default
closure route is absolute majorization.

\subsubsection{2.4 Typed KP activity is preferred over scalar
smallness}\label{typed-kp-activity-is-preferred-over-scalar-smallness}

Scalar KP screens are too crude. The 3D program should use a typed
activity vector

{[} \mathbf z(k,I\_k)=\left(z\_1(k,I\_k),\ldots,z\_T(k,I\_k)\right) {]}

with compatibility matrix (A\_\{ts\}\^{}\{3D\}(k)) and penalties
(a\_s(k)):

{[} \sum\emph{s A}\{ts\}\^{}\{3D\}(k)
\Bigl(z\_s\textsuperscript{\{\rm term\}(k,I\_k)+z\_s}\{\rm nonterm\}(k,I\_k)\Bigr)
e\^{}\{a\_s(k)\} \le a\_t(k). {]}

A scalar fallback condition such as

{[} 3\bigl(b\_\{3D\}(k,I\_k)+R\_\{3D\}(k,I\_k)\bigr)\textless1 {]}

may be used only as a conservative diagnostic or early prototype.

\subsubsection{2.5 A terminal bound should export a residual
budget}\label{a-terminal-bound-should-export-a-residual-budget}

A terminal graph is useful if it produces not only a pass/fail result
but also a residual budget for all remaining activity:

{[} R\_\{\max\}(k,I\_k) \quad\text{or}\quad \mathbf R\_\{\max\}(k,I\_k).
{]}

The full nonterminal program is then constrained by this exported
budget.

\subsubsection{2.6 Periodic volume uniformity is a separate theorem
fragment}\label{periodic-volume-uniformity-is-a-separate-theorem-fragment}

A finite local graph does not automatically imply a thermodynamic bound.
Every finite graph row must have a periodic-template successor proving a
volume-independent estimate:

{[} \sup\emph{\Lambda\sup}\{x\in\Lambda\}
\sum\emph{\{y\in\Lambda\}B}\Lambda(x,y;k,I\_k) \le b\_\{3D\}(k,I\_k).
{]}

\subsubsection{2.7 Source-audit failures are useful
failures}\label{source-audit-failures-are-useful-failures}

If a row fails because a required source object does not exist, the
architecture is working. In particular, no Taylor remainder bound may be
claimed until the nonterminal term inventory exists.

\subsubsection{2.8 Terminal success exports a vector, not an RG
rung}\label{terminal-success-exports-a-vector-not-an-rg-rung}

The terminal K-star pattern is a useful micro-lemma template, but it
must not be promoted into a full RG theorem. A terminal gate must export
a typed activity vector and an absorption contract, not merely a scalar
smallness number.

The required terminal output is a typed vector, with components for
link, plaquette, face, corner, cube, multiblock, gauge, and tail
activity.

The next-scale or full-RG row must prove an absorption statement

{[}
\mathbf z\_\{3D\}\textsuperscript{\{\rm term\}(k,I\_k)+\mathbf z\_\{3D\}}\{\rm nonterm\}(k,I\_k)
\in \mathcal D\_\{\rm conv\}(k,I\_k), {]}

where the declared convergence domain is KP, Dobrushin, or a tree-graph
improved polymer domain.

\begin{center}\rule{0.5\linewidth}{0.5pt}\end{center}

\subsection{3. Scale-Indexed Activity
Architecture}\label{scale-indexed-activity-architecture}

\subsubsection{3.1 Coupling slabs}\label{coupling-slabs}

At each RG scale (k), the coupling is represented by an interval slab

{[} I\_k={[}\beta\_k\^{}-,\beta\_k\^{}+{]}. {]}

All activity bounds must be uniform over (\beta\in I\_k).

\subsubsection{3.2 Terminal activity
weight}\label{terminal-activity-weight}

For a terminal atom class (\alpha) of type (t), define

{[} z\_\{t,\alpha\}\^{}\{\rm term\}(k,I\_k) = m\_\alpha
\sup\emph{\{\beta\in I\_k\}\textbar c}\alpha(\beta,k)\textbar{}
W\_\alpha(k), {]}

where:

\begin{itemize}
\tightlist
\item
  (m\_\alpha) is the exact orbit/multiplicity factor;
\item
  (c\_\alpha(\beta,k)) is the coefficient bound;
\item
  (W\_\alpha(k)) is the polymer weight, including support diameter,
  representation penalty, scale penalty, and locality penalty.
\end{itemize}

The full terminal activity type is

{[} z\_t\^{}\{\rm term\}(k,I\_k) =
\sum\emph{\{\alpha\in T\_t\}z}\{t,\alpha\}\^{}\{\rm term\}(k,I\_k). {]}

\subsubsection{3.3 Nonterminal activity
weight}\label{nonterminal-activity-weight}

The nonterminal activity is similarly decomposed:

{[} z\_t\^{}\{\rm nonterm\}(k,I\_k) = z\_t\^{}\{\rm Taylor\}
+z\_t\^{}\{\rm nonlinear\} +z\_t\^{}\{\rm multiblock\}
+z\_t\^{}\{\rm gauge/RP\} +z\_t\^{}\{\rm tail\}
+z\_t\^{}\{\rm leakage\}. {]}

Every summand must be generated by a source ledger and bounded uniformly
over (I\_k).

\subsubsection{3.4 Polymer convergence-domain
modularity}\label{polymer-convergence-domain-modularity}

The activity gate must not be hard-coded to the classical
Kotecky--Preiss criterion. The checker schema should allow a declared
convergence criterion

\begin{Shaded}
\begin{Highlighting}[]
\NormalTok{criterion ∈ \{KP, Dobrushin, Fernandez{-}Procacci, custom\_tree\_graph\}}
\end{Highlighting}
\end{Shaded}

with an accompanying proof row for the chosen sufficient condition. The
default criterion is KP, but the design must allow sharper tree-graph or
Penrose-identity bounds when absolute majorization is too crude.

For each criterion, the row must declare:

\begin{longtable}[]{@{}
  >{\raggedright\arraybackslash}p{(\columnwidth - 2\tabcolsep) * \real{0.5000}}
  >{\raggedright\arraybackslash}p{(\columnwidth - 2\tabcolsep) * \real{0.5000}}@{}}
\toprule\noalign{}
\begin{minipage}[b]{\linewidth}\raggedright
Field
\end{minipage} & \begin{minipage}[b]{\linewidth}\raggedright
Meaning
\end{minipage} \\
\midrule\noalign{}
\endhead
\bottomrule\noalign{}
\endlastfoot
\texttt{criterion\_id} & KP, Dobrushin, Fernandez--Procacci, or
custom \\
\texttt{activity\_vector\_id} & typed vector being tested \\
\texttt{compatibility\_matrix\_id} & graph or matrix used by the
criterion \\
\texttt{tree\_graph\_bound\_id} & required if using tree-graph
improvement \\
\texttt{domain\_certificate} & exact inequality defining the convergence
domain \\
\texttt{negative\_controls} & known overbudget vectors rejected by the
criterion \\
\end{longtable}

The fallback scalar gate remains diagnostic only.

\subsubsection{3.5 Dimensional-flow monotonicity with gauge-ripple
correction}\label{dimensional-flow-monotonicity-with-gauge-ripple-correction}

The scale ledger must prove that the RG flow moves toward the fallback
basin. If the effective map has the form

{[} \beta\_\{k+1\}=\frac{\beta_k}{L}+\Delta\_k, {]}

then the correction term must be decomposed, not treated as a single
unstructured error:

{[} \textbar{}\Delta\emph{k\textbar{}}\{\max\} =
\textbar{}\Delta\emph{k\^{}\{\rm analytic\}\textbar{}}\{\max\}
+\textbar{}\Delta\emph{k\^{}\{\rm trunc\}\textbar{}}\{\max\}
+\textbar{}\Delta\emph{k\^{}\{\rm gauge\}\textbar{}}\{\max\}
+\textbar{}\Delta\emph{k\^{}\{\rm tail\}\textbar{}}\{\max\}. {]}

Here ( \textbar{}\Delta\emph{k\^{}\{\rm gauge\}\textbar{}}\{\max\} )
records the analytic effect of changing local block-axial or tree-gauge
coordinates between scale (k) and scale (k+1). Such gauge-ripple terms
can cross block boundaries and must be included before a monotone
coupling descent is claimed.

The corrected monotonicity gate is

{[} \frac{\beta_k^+}{L}
+\textbar{}\Delta\emph{k\^{}\{\rm analytic\}\textbar{}}\{\max\}
+\textbar{}\Delta\emph{k\^{}\{\rm trunc\}\textbar{}}\{\max\}
+\textbar{}\Delta\emph{k\^{}\{\rm gauge\}\textbar{}}\{\max\}
+\textbar{}\Delta\emph{k\^{}\{\rm tail\}\textbar{}}\{\max\}
\textless{}\beta\_k\^{}-. {]}

Equivalently, in interval form,

{[} I\_\{k+1\}\textsuperscript{+\textless I\_k}-. {]}

All four components of (\textbar{}\Delta\emph{k\textbar{}}\{\max\}) must
be computed in the same declared metric space. The scale-coupling ledger
must include a \texttt{flow\_error\_norm} object with fields

\begin{longtable}[]{@{}
  >{\raggedright\arraybackslash}p{(\columnwidth - 2\tabcolsep) * \real{0.5000}}
  >{\raggedright\arraybackslash}p{(\columnwidth - 2\tabcolsep) * \real{0.5000}}@{}}
\toprule\noalign{}
\begin{minipage}[b]{\linewidth}\raggedright
Field
\end{minipage} & \begin{minipage}[b]{\linewidth}\raggedright
Meaning
\end{minipage} \\
\midrule\noalign{}
\endhead
\bottomrule\noalign{}
\endlastfoot
\texttt{norm\_id} & stable identifier for the norm used by all
(\Delta\_k) components \\
\texttt{space} & \texttt{effective\_coupling\_parameter},
\texttt{effective\_action\_density}, \texttt{polymer\_activity\_norm},
or \texttt{transfer\_operator\_norm} \\
\texttt{domain} & scale slab and volume class over which the supremum is
taken \\
\texttt{uniform\_over\_beta} & whether the estimate is uniform over
(eta\in I\_k) \\
\texttt{uniform\_over\_volume} & whether the estimate is volume
independent \\
\texttt{unit} & units used for comparison with (eta) \\
\texttt{conversion\_certificate\_ids} & certificates converting any
component proved in another norm \\
\end{longtable}

If a component is naturally bounded in another norm, the row must
include a conversion certificate

{[} \textbar{}\Delta\textbar{}\emph{\{eta\}\le C}\{ m
conv\}\textbar{}\Delta\textbar\_\{\mathcal N\}. {]}

No dimensional-flow row may add heterogeneous error components unless
all have been converted into the declared \texttt{flow\_error\_norm}.

Any exception must be explicitly declared and bounded. The default
production blocking factor is dyadic,

{[} L=2. {]}

If the inequality fails at (L=2), the standard remedies are to sharpen
field-localization bounds, reduce or rescale the small-field threshold
(p\_k), refine the scale slab (I\_k), improve the gauge-ripple estimate,
or split the block phase cell more finely. Blindly increasing (L) is not
an accepted default remedy, because the block volume grows as (L\^{}3),
increasing internal phase space and potentially enlarging
(\textbar{}\Delta\emph{k\^{}\{\rm gauge\}\textbar{}}\{\max\}). Larger
blocking factors may be explored only in an explicitly marked
experimental lane with their own gauge-ripple and boundary-fluctuation
budget.

\begin{center}\rule{0.5\linewidth}{0.5pt}\end{center}

\subsection{4. Gauge Handling and Reflection-Positivity
Discipline}\label{gauge-handling-and-reflection-positivity-discipline}

\subsubsection{4.1 Gauge lanes}\label{gauge-lanes}

Every local atom, term, or polymer activity must carry a
\texttt{gauge\_lane} field:

\begin{longtable}[]{@{}
  >{\raggedright\arraybackslash}p{(\columnwidth - 2\tabcolsep) * \real{0.5000}}
  >{\raggedright\arraybackslash}p{(\columnwidth - 2\tabcolsep) * \real{0.5000}}@{}}
\toprule\noalign{}
\begin{minipage}[b]{\linewidth}\raggedright
Lane
\end{minipage} & \begin{minipage}[b]{\linewidth}\raggedright
Meaning
\end{minipage} \\
\midrule\noalign{}
\endhead
\bottomrule\noalign{}
\endlastfoot
\texttt{invariant} & Gauge-invariant object, safe for external
observables if also RP-compatible \\
\texttt{block\_axial} & Object expressed in a block axial/tree gauge \\
\texttt{tree\_gauge} & Object expressed after local spanning-tree gauge
fixing \\
\texttt{residual} & Object depending on residual boundary gauge
variables \\
\texttt{gauge\_covariant\_background} & Object aligned to a covariant
background description for IR/fallback use \\
\texttt{averaged} & Object transferred back by explicit gauge
averaging \\
\texttt{internal\_only} & Proof-device object not allowed in final
OS-facing claims \\
\end{longtable}

A gauge-lane assignment alone is not sufficient. If an object enters a
covariant-background lane or any lane with a nontrivial Jacobian, the
atom or term must link to a determinant certificate rather than assuming
the tree-gauge Haar identity. The base atom schema therefore includes
optional hooks \texttt{jacobian\_type},
\texttt{jacobian\_determinant\_id}, and
\texttt{ghost\_field\_dependence}.

\subsubsection{4.2 Tree-axial gauge
metadata}\label{tree-axial-gauge-metadata}

A Balaban-style maximally tree-axial gauge option is included. For each
block (B), a spanning tree (T\_B) is declared and links on the tree are
fixed:

{[} U\_b=\mathbf 1, \qquad b\in T\_B. {]}

The corresponding ledger must verify:

\begin{enumerate}
\def\labelenumi{\arabic{enumi}.}
\tightlist
\item
  (T\_B) is connected;
\item
  (T\_B) contains no cycles;
\item
  (T\_B) touches all block vertices;
\item
  every interior link has bounded distance to the block boundary;
\item
  residual gauge transformations are represented exactly;
\item
  the Haar measure change of variables is exactly ledgered.
\end{enumerate}

The lanes \texttt{tree\_gauge} and \texttt{block\_axial} are local
UV/intermediate-scale coordinate lanes only. They are not permanent
infrared descriptions. Repeated block transformations must eventually
transfer these objects through a certified gauge-covariant or
gauge-invariant lane before fallback-basin or clustering claims are
made. The required long-scale transition has the schematic form

\begin{Shaded}
\begin{Highlighting}[]
\NormalTok{tree\_gauge or block\_axial}
\NormalTok{  {-}\textgreater{} gauge\_covariant\_background}
\NormalTok{  {-}\textgreater{} haar\_averaged\_invariant}
\NormalTok{  {-}\textgreater{} polymer\_or\_character\_activity}
\end{Highlighting}
\end{Shaded}

A row that leaves tree/block-axial discontinuities to propagate
indefinitely into the infrared fails, because block-boundary coordinate
jumps can generate nonlocal artifacts in effective polymer weights.

\subsubsection{4.3 Haar reversibility}\label{haar-reversibility}

For a gauge-fixing map

{[} U\mapsto (U\^{}\{\rm gf\},h), {]}

the checker must verify

{[} dU=dU\^{}\{\rm gf\},dh,J\_\{\rm gf\}. {]}

For compact finite-dimensional lattice tree gauges, the expected
Jacobian may be (J\_\{\rm gf\}=1), but this must be a proved ledger row.
The residual Haar integral must satisfy

{[} \int dh=1 {]}

with normalized Haar measure.

\subsubsection{4.4 Configuration-dependent determinant
hooks}\label{configuration-dependent-determinant-hooks}

Tree and block-axial gauges in finite-dimensional lattice variables may
have a constant Haar Jacobian, often certified as

{[} J\_\{\rm gf\}=1. {]}

This simplification is restricted to the declared tree/block-axial lane.
If repeated block transformations transfer an object into a
gauge-covariant background lane, the row must allow a
configuration-dependent Faddeev--Popov or determinant factor. Such a
factor is not a static coefficient; it may depend on background and
fluctuation variables and may generate nonterminal ghost or
determinant-loop terms.

The determinant metadata is represented by optional fields:

\begin{longtable}[]{@{}
  >{\raggedright\arraybackslash}p{(\columnwidth - 2\tabcolsep) * \real{0.5000}}
  >{\raggedright\arraybackslash}p{(\columnwidth - 2\tabcolsep) * \real{0.5000}}@{}}
\toprule\noalign{}
\begin{minipage}[b]{\linewidth}\raggedright
Field
\end{minipage} & \begin{minipage}[b]{\linewidth}\raggedright
Meaning
\end{minipage} \\
\midrule\noalign{}
\endhead
\bottomrule\noalign{}
\endlastfoot
\texttt{jacobian\_type} & \texttt{constant\_haar},
\texttt{configuration\_dependent\_fp}, \texttt{ghost\_expanded}, or
\texttt{determinant\_ratio} \\
\texttt{jacobian\_determinant\_id} & link to determinant certificate \\
\texttt{ghost\_field\_dependence} & whether ghost variables or ghost
loops appear \\
\texttt{determinant\_expansion\_id} & expansion ledger for nontrivial
determinant terms \\
\texttt{ghost\_loop\_degree} & loop/degree order if ghost expansion is
used \\
\texttt{background\_field\_id} & covariant-background field reference \\
\texttt{fp\_operator\_id} & discrete FP operator or determinant matrix
reference \\
\end{longtable}

If \texttt{gauge\_lane\ =\ gauge\_covariant\_background}, then
determinant metadata is mandatory unless a row proves that the
determinant contribution cancels or is identically one in the declared
finite setting.

A determinant/ghost expansion is not automatically admissible in a
positive polymer criterion. Grassmann variables do not define a positive
measure in the ordinary bosonic sense, and ghost loops carry signs.
Therefore any ghost-expanded determinant contribution must pass

\begin{Shaded}
\begin{Highlighting}[]
\NormalTok{YM3D\_GHOST\_DETERMINANT\_SIGN\_CONSISTENCY\_CHECK\_1}
\end{Highlighting}
\end{Shaded}

before it can enter a convergence gate. The row must distinguish among
exact determinant retention, ghost expansion, determinant ratios, and
perturbative determinant truncations. If
\texttt{ghost\_field\_dependence=true}, the row must specify whether the
contribution is handled by a signed/fermionic cluster expansion, by an
absolute-majorant certificate, or by an exact determinant-positivity
certificate. A ghost-dependent row may not enter an ordinary positive KP
activity vector unless \texttt{positive\_activity\_admissible=true} is
imported from the determinant sign certificate.

The determinant sign certificate records:

\begin{longtable}[]{@{}
  >{\raggedright\arraybackslash}p{(\columnwidth - 2\tabcolsep) * \real{0.5000}}
  >{\raggedright\arraybackslash}p{(\columnwidth - 2\tabcolsep) * \real{0.5000}}@{}}
\toprule\noalign{}
\begin{minipage}[b]{\linewidth}\raggedright
Field
\end{minipage} & \begin{minipage}[b]{\linewidth}\raggedright
Meaning
\end{minipage} \\
\midrule\noalign{}
\endhead
\bottomrule\noalign{}
\endlastfoot
\texttt{determinant\_sign\_certificate\_id} & sign/positivity or
signed-cluster certificate \\
\texttt{requires\_signed\_cluster\_certificate} & whether ordinary
positive activity criteria are forbidden \\
\texttt{positive\_activity\_admissible} & whether the term may enter a
positive polymer bound \\
\texttt{absolute\_majorant\_certificate\_id} & certificate replacing
signed ghost loops by a positive majorant \\
\texttt{determinant\_ratio\_lower\_bound\_id} & lower-bound certificate
for determinant denominators \\
\texttt{ghost\_loop\_sign\_rule} & convention for loop signs and
anticommutation ordering \\
\end{longtable}

\subsubsection{4.5 Gauge-ripple flow
correction}\label{gauge-ripple-flow-correction}

Because gauge fixing is performed locally inside coarse blocks, changing
the block scale changes the block trees, residual boundary variables,
and gauge-averaging maps. The difference between the scale-(k) and
scale-(k+1) gauge coordinates is a proof-device effect, but its analytic
size can affect the certified coupling-flow interval.

The ledger \texttt{YM3D\_GAUGE\_RIPPLE\_FLOW\_CORRECTION\_1} records
this effect through

{[} \textbar{}\Delta\emph{k\^{}\{\rm gauge\}\textbar{}}\{\max\}. {]}

Each row must include:

\begin{longtable}[]{@{}
  >{\raggedright\arraybackslash}p{(\columnwidth - 2\tabcolsep) * \real{0.5000}}
  >{\raggedright\arraybackslash}p{(\columnwidth - 2\tabcolsep) * \real{0.5000}}@{}}
\toprule\noalign{}
\begin{minipage}[b]{\linewidth}\raggedright
Field
\end{minipage} & \begin{minipage}[b]{\linewidth}\raggedright
Meaning
\end{minipage} \\
\midrule\noalign{}
\endhead
\bottomrule\noalign{}
\endlastfoot
\texttt{scale\_slab\_id} & source interval (I\_k) \\
\texttt{block\_size\_L} & blocking factor \\
\texttt{source\_tree\_id} & scale-(k) tree/axial gauge data \\
\texttt{target\_tree\_id} & scale-(k+1) tree/axial gauge data \\
\texttt{boundary\_orbit\_id} & residual block-boundary gauge orbit \\
\texttt{ripple\_support} & support affected by the gauge-coordinate
transition \\
\texttt{analytic\_bound\_source} & row proving the ripple estimate \\
\texttt{delta\_gauge\_bound} & certified value of ( \\
\texttt{flow\_gate\_link} & monotonicity row using this term \\
\end{longtable}

No dimensional-flow monotonicity row may pass unless the corresponding
gauge-ripple contribution is either zero by exact symmetry or included
explicitly in the (\Delta\_k) budget.

\subsubsection{4.6 Reflection-positivity
discipline}\label{reflection-positivity-discipline}

The Wilson lattice measure is the OS-positive object. Gauge-fixed,
block-averaged, or internally smeared variables are proof devices unless
transferred to gauge-invariant, reflection-compatible objects.

No OS-facing claim may be made from a gauge-fixed row unless the row has
passed:

\begin{enumerate}
\def\labelenumi{\arabic{enumi}.}
\tightlist
\item
  gauge averaging;
\item
  Haar reversibility;
\item
  reflection-parity tagging;
\item
  reflection-plane factorization;
\item
  an explicit \texttt{rp\_safe\_use} check.
\end{enumerate}

The reflection-plane factorization condition is geometric. For each
reflection plane (\Pi) used in an OS-facing claim, the Haar-averaging
support must either be contained in one half-space or be exactly
symmetric under reflection:

{[} H\_\{\rm avg\}\subset \Pi\^{}+, \qquad
H\_\{\rm avg\}\subset \Pi\^{}-, \qquad \text{or} \qquad
\theta(H\_\{\rm avg\})=H\_\{\rm avg\}. {]}

A support that crosses the reflection plane asymmetrically fails Gate
S12 unless a separate exact factorization proof is attached. The
corresponding certificate records
\texttt{spatial\_factorization\_plane}, \texttt{support\_relation}, and
\texttt{factorization\_certificate\_id}.

\begin{center}\rule{0.5\linewidth}{0.5pt}\end{center}

\subsection{5. Small/Large-Field Partition
Policy}\label{smalllarge-field-partition-policy}

\subsubsection{5.1 Default choice: sharp characteristic
cutoffs}\label{default-choice-sharp-characteristic-cutoffs}

The first implementation uses sharp cutoffs:

{[} \chi\emph{\{\rm SF\}(A)+\chi}\{\rm LF\}(A)=1. {]}

The small-field region may be defined, for example, by

{[} \max\_\{p\subset B\} d(U\_p,\mathbf 1)\le p\_k, {]}

or by a gauge-fixed block curvature norm

{[} \textbar F\_A\textbar\_\{B,k\}\le p\_k. {]}

\subsubsection{5.2 Sharp-cutoff rule}\label{sharp-cutoff-rule}

No Taylor, BKAR, cluster, or interpolation derivative may hit a sharp
cutoff unless a boundary-integration subledger is supplied.

If a derivative hits (\chi\emph{\{\rm SF\}) or (\chi}\{\rm LF\}), the
row fails unless the resulting boundary contribution is explicitly
integrated and placed in a valid residual bucket.

A valid sharp-boundary row must include a boundary-surface multiplier
certificate. The bound has the form

{[} \textbar{}\mathcal B\textbar{} \le
\textbar f\textbar{}\emph{\{\infty\} ,M}\{\partial{\rm SF}\}, {]}

where (M\_\{\partial{\rm SF}\}) is an exact rational or interval
enclosure of the relevant configuration-space hypersurface volume,
coarea factor, or trace-theorem constant. The nonterminal row must
record \texttt{boundary\_surface\_multiplier\_id},
\texttt{trace\_operator\_id}, and
\texttt{boundary\_measure\_enclosure\_id} whenever a sharp cutoff
boundary contributes.

\subsubsection{5.3 Smooth-cutoff upgrade
branch}\label{smooth-cutoff-upgrade-branch}

A later smooth branch may use Gevrey-class partitions

{[} \chi\emph{\{\rm core\}+\chi}\{\rm trans\}+\chi\_\{\rm tail\}=1 {]}

with derivative bounds

{[} \textbar{}\partial\^{}m\chi\textbar{}
\le C\_\chi\textsuperscript{\{m+1\}(m!)}\sigma p\_k\^{}\{-m\}. {]}

The derivative leakage gate must then prove that these losses fit inside
a declared leakage budget:

{[} R\_\{\rm leakage\}(k,I\_k) \le R\_\{\rm leakage,max\}(k,I\_k). {]}

\begin{center}\rule{0.5\linewidth}{0.5pt}\end{center}

\subsection{6. Symmetry-Reduced Local
Geometry}\label{symmetry-reduced-local-geometry}

\subsubsection{6.1 Cubic group representation and orbit
representatives}\label{cubic-group-representation-and-orbit-representatives}

The production generator represents the full local cubic group (O\_h) by
signed permutation matrices

{[} M\in \mathrm{Mat}\_\{3\times 3\}(\mathbb Z), \qquad M\^{}TM=I,
\qquad M e\_i=\pm e\_j. {]}

Unless a row explicitly restricts to the orientation-preserving
subgroup, the local symmetry group has

{[} \textbar O\_h\textbar=48. {]}

This signed-matrix representation is the source of truth. It induces
permutations of oriented links, plaquettes, supports, anchors, and
collision pairs. For an oriented link (\ell=(x,\mu)) from (x) to
(x+e\_\mu), if

{[} M e\_\mu=s e\_\nu, \qquad s\in\{+1,-1\}, {]}

then for (s=+1),

{[} M\cdot(x,\mu)=(Mx,\nu), {]}

while for (s=-1), the transformed negative-oriented link is rewritten as
a positive-oriented link with an orientation-flip tag:

{[} M\cdot(x,\mu)=(Mx-e\_\nu,\nu). {]}

The matrix action on links lifts upward to all composite supports by a
mandatory composition rule:

\begin{enumerate}
\def\labelenumi{\arabic{enumi}.}
\tightlist
\item
  apply the signed permutation matrix to every constituent oriented
  link;
\item
  canonicalize each transformed link to the positive-axis convention
  with an orientation-flip tag;
\item
  sort the transformed links into a canonical integer tuple;
\item
  compute the transformed support identifier from this tuple;
\item
  verify that plaquette, cube, and multiblock transforms agree with the
  transforms of their constituent links.
\end{enumerate}

The orbit schema must record:

\begin{longtable}[]{@{}
  >{\raggedright\arraybackslash}p{(\columnwidth - 2\tabcolsep) * \real{0.5000}}
  >{\raggedright\arraybackslash}p{(\columnwidth - 2\tabcolsep) * \real{0.5000}}@{}}
\toprule\noalign{}
\begin{minipage}[b]{\linewidth}\raggedright
Field
\end{minipage} & \begin{minipage}[b]{\linewidth}\raggedright
Meaning
\end{minipage} \\
\midrule\noalign{}
\endhead
\bottomrule\noalign{}
\endlastfoot
\texttt{group\_action\_source} & \texttt{signed\_permutation\_matrix} \\
\texttt{fundamental\_action} & \texttt{oriented\_link\_action} \\
\texttt{composition\_rule\_id} & canonical link-tuple lift rule \\
\texttt{support\_lift\_rule.plaquette} & ordered boundary-link lift \\
\texttt{support\_lift\_rule.cube} & face or boundary-link lift \\
\texttt{support\_lift\_rule.multiblock} & canonical union of block
supports \\
\texttt{orientation\_canonicalization} & positive-axis plus flip tag \\
\texttt{canonical\_hash\_rule} & deterministic lexicographic integer
tuple rule \\
\end{longtable}

An independent audit mode reconstructs the same orbit partition by pure
link-coordinate transformations and verifies equality with the
signed-matrix implementation.

The local cubic symmetry group acts on atoms, supports, anchors, and
collision pairs. The program stores orbit representatives ({[}\alpha{]})
with exact stabilizers:

{[}
m\_\{{[}\alpha{]}\}=\frac{|G_{\rm local}|}{|\operatorname{Stab}(\alpha)|}.
{]}

For the default full cubic group, this specializes to

{[}
\textbar{}\operatorname{Orb}(\alpha)\textbar=\frac{48}{|\operatorname{Stab}(\alpha)|}.
{]}

Collision edges are likewise represented by orbit pairs

{[} ({[}\alpha{]},{[}\beta{]},\mathrm{collision\_type}) {]}

with exact multiplicity.

\paragraph{Orbit-closure consistency}\label{orbit-closure-consistency}

The orbit schema must also include an explicit orbit-closure proof. For
each representative ({[}\alpha{]}), with canonical representative map
(\operatorname{rep}), the checker verifies

{[} \forall g\in O\_h,\qquad
\operatorname{rep}(g\cdot r\_\alpha)=r\_\alpha. {]}

Equivalently, for every concrete atom (\alpha) and every group element
(g),

{[} \operatorname{rep}(g\cdot \alpha)=\operatorname{rep}(\alpha) {]}

whenever (g\cdot\alpha) lies in the same orbit. The representative map
must also be idempotent:

{[} \operatorname{rep}(\operatorname{rep}(\alpha))=
\operatorname{rep}(\alpha). {]}

The checker row

\begin{Shaded}
\begin{Highlighting}[]
\NormalTok{YM3D\_ORBIT\_CLOSURE\_CONSISTENCY\_CHECK\_1}
\end{Highlighting}
\end{Shaded}

is a hard prerequisite for the expanded-graph audit. It ensures that
inconsistent choices of representatives at the link, plaquette, cube, or
multiblock level cannot corrupt stabilizer sizes or majorant
multiplicities.

\subsubsection{6.2 Orbit-reduced majorant}\label{orbit-reduced-majorant}

The typed majorant matrix is computed from orbit data:

{[} B\_\{ts\}\^{}\{3D\}(k,I\_k) = \sum\emph{\{{[}\alpha{]}\in T\_t\}
\sum}\{{[}\beta{]}\in T\_s\} m\_\{{[}\alpha{]},{[}\beta{]}\}
w\_\{{[}\alpha{]},{[}\beta{]}\}(k,I\_k). {]}

The full expanded graph is generated only for audit boxes and negative
controls.

\subsubsection{6.3 Expanded-graph audit on small
lattices}\label{expanded-graph-audit-on-small-lattices}

Orbit reduction is a production representation, not a substitute for
validation. The validator must generate fully expanded collision graphs
on designated small periodic audit lattices and compare them to the
orbit-reduced graph by explicit projection and multiplicity checks.

Let (G\_\{\rm exp\}=(V\_\{\rm exp\},E\_\{\rm exp\})) be the fully
expanded atom collision graph, and let

{[} \pi:V\_\{\rm exp\}\to V\_\{\rm orb\} {]}

be the projection to orbit representatives. For every orbit-pair
collision type (q), the checker must verify

{[} M\^{}\{\rm expanded\}\emph{\{{[}\alpha{]},{[}\beta{]},q\} =
\#\{(u,v)\in E}\{\rm exp\}:\pi(u)={[}\alpha{]},\pi(v)={[}\beta{]},~q(u,v)=q\}
= M\^{}\{\rm reduced\}\_\{{[}\alpha{]},{[}\beta{]},q\}. {]}

It must also verify stabilizer multiplicities

{[}
m\_\{{[}\alpha{]}\}=\frac{|G_{\rm local}|}{|\operatorname{Stab}(\alpha)|}.
{]}

This audit is a hard acceptance gate for
\texttt{YM3D\_CUBIC\_SYMMETRY\_ORBIT\_REDUCTION\_1}. If the expanded
graph and reduced graph disagree on any audit box, the orbit-reduction
row fails.

\begin{center}\rule{0.5\linewidth}{0.5pt}\end{center}

\subsection{7. K-Star Terminal Graph
Program}\label{k-star-terminal-graph-program}

\subsubsection{7.1 Terminal atom export}\label{terminal-atom-export}

The first 3D K-star artifact is

\texttt{YM3D\_KSTAR\_TERMINAL\_ATOM\_EXPORT\_1}.

Each atom row must include:

\begin{longtable}[]{@{}
  >{\raggedright\arraybackslash}p{(\columnwidth - 2\tabcolsep) * \real{0.5000}}
  >{\raggedright\arraybackslash}p{(\columnwidth - 2\tabcolsep) * \real{0.5000}}@{}}
\toprule\noalign{}
\begin{minipage}[b]{\linewidth}\raggedright
Field
\end{minipage} & \begin{minipage}[b]{\linewidth}\raggedright
Meaning
\end{minipage} \\
\midrule\noalign{}
\endhead
\bottomrule\noalign{}
\endlastfoot
\texttt{atom\_id} & stable unique identifier \\
\texttt{orbit\_id} & cubic symmetry orbit \\
\texttt{support\_id} & local support \\
\texttt{anchor\_type} & link, plaquette, cube, face, edge, corner,
multiblock \\
\texttt{orientation} & declared orientation convention \\
\texttt{representation\_label} & spin/character/spin-network label \\
\texttt{gauge\_lane} & gauge state of the object \\
\texttt{rp\_parity\_tag} & reflection behavior \\
\texttt{coefficient\_formula} & exact symbolic coefficient \\
\texttt{coefficient\_bound} & interval/rational upper bound over
(I\_k) \\
\texttt{scale\_slab} & (k,I\_k) validity \\
\end{longtable}

\subsubsection{7.2 Collision graph}\label{collision-graph}

The terminal collision graph records geometric overlap, shared boundary
variables, shared gauge residues, or support incompatibility. Every edge
carries a typed collision label.

\subsubsection{7.3 Terminal KP gate}\label{terminal-kp-gate}

The terminal-only gate is

{[} \sum\emph{s
A}\{ts\}\textsuperscript{\{3D\}(k)z\_s}\{\rm term\}(k,I\_k)e\^{}\{a\_s(k)\}
\le a\_t(k). {]}

Passing this gate exports a residual budget for nonterminal activity. It
does not close the full RG rung.

\subsubsection{7.4 Terminal vector-export
contract}\label{terminal-vector-export-contract}

A terminal row must export the full vector

{[} \mathbf z\_\{3D\}\^{}\{\rm term\}(k,I\_k) {]}

plus a machine-readable contract describing which nonterminal types may
absorb, refine, or interact with each component. The contract must
include:

\begin{longtable}[]{@{}
  >{\raggedright\arraybackslash}p{(\columnwidth - 2\tabcolsep) * \real{0.5000}}
  >{\raggedright\arraybackslash}p{(\columnwidth - 2\tabcolsep) * \real{0.5000}}@{}}
\toprule\noalign{}
\begin{minipage}[b]{\linewidth}\raggedright
Field
\end{minipage} & \begin{minipage}[b]{\linewidth}\raggedright
Meaning
\end{minipage} \\
\midrule\noalign{}
\endhead
\bottomrule\noalign{}
\endlastfoot
\texttt{terminal\_vector\_id} & identifier for the terminal activity
vector \\
\texttt{component\_type} & link, plaquette, face, corner, cube,
multiblock, gauge, tail \\
\texttt{bound\_value} & exact rational/interval bound for that
component \\
\texttt{scale\_slab\_id} & scale/coupling interval \\
\texttt{absorption\_target} & nonterminal bucket or next-scale object
that absorbs it \\
\texttt{convergence\_domain\_id} &
KP/Dobrushin/Fernandez--Procacci/custom domain \\
\texttt{nonclaim} & explicit statement that terminal success is not an
RG rung \\
\end{longtable}

A scalar terminal bound may be recorded for readability, but it is not
sufficient for a full RG transition.

\begin{center}\rule{0.5\linewidth}{0.5pt}\end{center}

\subsection{8. Nonterminal Taylor and Forest Remainder
Program}\label{nonterminal-taylor-and-forest-remainder-program}

\subsubsection{8.1 Required inventory
fields}\label{required-inventory-fields}

Every nonterminal Taylor or forest-generated term must include:

\begin{longtable}[]{@{}
  >{\raggedright\arraybackslash}p{(\columnwidth - 2\tabcolsep) * \real{0.5000}}
  >{\raggedright\arraybackslash}p{(\columnwidth - 2\tabcolsep) * \real{0.5000}}@{}}
\toprule\noalign{}
\begin{minipage}[b]{\linewidth}\raggedright
Field
\end{minipage} & \begin{minipage}[b]{\linewidth}\raggedright
Meaning
\end{minipage} \\
\midrule\noalign{}
\endhead
\bottomrule\noalign{}
\endlastfoot
\texttt{term\_id} & unique term identifier \\
\texttt{source\_expansion} & Taylor, cumulant, BKAR, character, or other
source \\
\texttt{forest\_id} & forest identifier if BKAR/interpolation is used \\
\texttt{support\_id} & local support \\
\texttt{anchor\_type} & support anchor \\
\texttt{representation\_label} & spin/character/spin-network label \\
\texttt{degree} & Taylor or derivative degree \\
\texttt{coefficient\_formula} & exact symbolic coefficient \\
\texttt{interval\_upper\_bound} & certified upper bound over (I\_k) \\
\texttt{multiplicity} & exact local multiplicity \\
\texttt{bucket} & Taylor, nonlinear, multiblock, gauge/RP, tail,
leakage \\
\texttt{small\_field\_cutoff\_id} & valid Taylor domain \\
\texttt{field\_norm} & norm used for small-field region \\
\texttt{cutoff\_threshold} & (p\_k) or equivalent \\
\texttt{boundary\_leakage\_bound} & if applicable \\
\texttt{boundary\_surface\_multiplier\_id} & boundary/trace certificate
if a sharp cutoff boundary is used \\
\texttt{primitive\_enclosure\_ids} & analytic primitive intervals
imported by the term \\
\texttt{jacobian\_determinant\_id} & determinant/FP certificate if the
term contains nontrivial Jacobian data \\
\texttt{determinant\_expansion\_id} & expansion ledger for
configuration-dependent determinants \\
\texttt{ghost\_field\_dependence} & whether the term depends on ghost
variables \\
\texttt{ghost\_loop\_degree} & degree/order of ghost loop if
applicable \\
\texttt{ghost\_measure\_id} & Berezin/ghost integration measure
certificate if ghost fields appear \\
\texttt{ghost\_cutoff\_policy} & \texttt{none},
\texttt{inherited\_from\_gauge}, or \texttt{separate\_grassmann} \\
\texttt{grassmann\_cutoff\_id} & optional algebraic/Grassmann cutoff or
truncation certificate \\
\texttt{tail\_bucket\_link} & complement/tail assignment \\
\texttt{parent\_rows} & required source rows \\
\end{longtable}

Configuration-dependent determinant terms must never be compressed into
an untyped coefficient. If \texttt{determinant\_expansion\_id} is
present, the row must also specify the determinant source,
background-field dependence, fluctuation-field dependence, and the
certificate proving the interval enclosure for each
determinant-generated primitive.

If \texttt{ghost\_field\_dependence=true}, the row must provide a
\texttt{ghost\_measure\_id}. Grassmann variables are not ordinary
bosonic amplitude variables; therefore the design uses the neutral field
\texttt{ghost\_measure\_id} and an explicit
\texttt{ghost\_cutoff\_policy} rather than assuming that ghost fields
inherit the gauge-field small/large-field partition. A separate
\texttt{grassmann\_cutoff\_id} is required only if a later analytic lane
introduces an algebraic ghost truncation or cutoff boundary.

\subsubsection{8.2 BKAR/forest
declaration}\label{bkarforest-declaration}

If BKAR or another forest formula is used, the inventory must declare:

\begin{itemize}
\tightlist
\item
  forest edge set;
\item
  interpolation parameters;
\item
  differentiated factors;
\item
  induced support;
\item
  connectivity class;
\item
  symmetry factor;
\item
  derivative order;
\item
  cutoff-interaction status.
\end{itemize}

The multiblock combinatorics must match the typed compatibility matrix.

\begin{center}\rule{0.5\linewidth}{0.5pt}\end{center}

\subsection{9. Large-Field and Defect-Tail
Ledger}\label{large-field-and-defect-tail-ledger}

The large-field complement is represented as thermodynamic defects. The
ledger must not assume that all non-abelian defects are closed dual
surfaces measured only by area. The schema supports a vector-valued
support measure

{[}
\mathbf m(D)=(A(D),\ell(D),V(D),n\_\partial(D),J(D),B(D),E(D),\ldots),
{]}

where (A) is surface area, (\ell) is line length or tube length, (V) is
localized block volume, (n\_\partial) records boundary-pinned
contributions, (J) records junctions, (B) records branching nodes, and
(E) records endpoints or attachment points when relevant. This allows
the tail ledger to represent co-dimension-one center-vortex-like sheets,
line-like tubes, localized high-curvature lumps, boundary-pinned
defects, and mixed complexes in which lower-dimensional and
higher-dimensional supports intersect.

Each defect row includes:

\begin{longtable}[]{@{}
  >{\raggedright\arraybackslash}p{(\columnwidth - 2\tabcolsep) * \real{0.5000}}
  >{\raggedright\arraybackslash}p{(\columnwidth - 2\tabcolsep) * \real{0.5000}}@{}}
\toprule\noalign{}
\begin{minipage}[b]{\linewidth}\raggedright
Field
\end{minipage} & \begin{minipage}[b]{\linewidth}\raggedright
Meaning
\end{minipage} \\
\midrule\noalign{}
\endhead
\bottomrule\noalign{}
\endlastfoot
\texttt{defect\_id} & unique defect identifier \\
\texttt{support} & support set \\
\texttt{scale} & RG scale (k) \\
\texttt{scale\_slab\_id} & coupling interval (I\_k) \\
\texttt{defect\_type} & surface, tube, localized\_high\_curvature,
boundary\_pinned, mixed, gauge, tail \\
\texttt{support\_codimension} & declared codimension or mixed-complex
label \\
\texttt{support\_measure\_vector} & vector
(\mathbf m=(A,\ell,V,n\_\partial,J,B,E,\ldots)) \\
\texttt{area\_component} & surface/plaquette area, if applicable \\
\texttt{length\_component} & line/tube length, if applicable \\
\texttt{volume\_component} & localized block volume, if applicable \\
\texttt{junction\_count} & number of topology-changing junctions \\
\texttt{branching\_node\_count} & number of branching nodes in mixed
complexes \\
\texttt{endpoint\_count} & number of endpoints or boundary
attachments \\
\texttt{topological\_modifier\_rate} & entropy-rate component assigned
to junction/branching data \\
\texttt{defect\_topology\_type} & closed surface, vortex sheet, tube,
lump, boundary-pinned defect, or mixed complex \\
\texttt{action\_cost\_vector} & coefficient vector for scale-slab action
suppression \\
\texttt{entropy\_model\_id} & source row for the entropy bound \\
\texttt{entropy\_rate\_vector} & entropy-rate vector for the declared
support measure \\
\texttt{entropy\_prefactor} & prefactor (C\_\{\rm ent\}) \\
\texttt{dominance\_margin} & certified lower bound for action-cost minus
entropy-rate against the support measure \\
\texttt{tail\_sum\_bound} & resulting geometric/exponential sum \\
\texttt{suppression\_factor} & exponential or character-tail
suppression \\
\texttt{interaction\_range} & compatibility range \\
\texttt{kp\_type} & typed activity class \\
\end{longtable}

The tail gate checks estimates of the form

{[} \sum\emph{\{D\ni x\} M(D)
\exp\left[-\mathbf c_{\rm tail}(k,I_k)\cdot \mathbf m(D)\right]\le R}\{\rm tail\}(k,I\_k).
{]}

\subsubsection{9.1 Defect entropy
dominance}\label{defect-entropy-dominance}

For each declared defect class, the ledger must provide an entropy
estimate

{[} N(\mathbf m)\le C\_\{\rm ent\}
\exp\left(\boldsymbol\mu\_\{\rm ent\}\cdot \mathbf m\right). {]}

The large-field tail row may pass only if action cost strictly dominates
entropy over the full scale slab. In vector form, the required condition
is

{[}
\left(\mathbf c\_\{\rm tail\}(k,I\_k)-\boldsymbol\mu\emph{\{\rm ent\}\right)\cdot\mathbf m
\ge \delta}\{\rm tail\}\textbar{}\mathbf m\textbar{} {]}

for all admissible support-measure vectors (\mathbf m), with
(\delta\_\{\rm tail\}\textgreater0) recorded. For a pure surface-defect
specialization, this reduces to

{[}
c\_\{\rm tail\}(k,I\_k)-\mu\emph{\{\rm ent\}\ge\delta}\{\rm tail\}\textgreater0.
{]}

This generalization is required because the Göpfert--Mack compact (U(1))
surface-defect template is only a structural guide. The non-abelian
(SU(2)) ledger must also audit lower-dimensional or localized
high-curvature defect supports.

Required artifact:

\begin{Shaded}
\begin{Highlighting}[]
\NormalTok{YM3D\_DEFECT\_SUPPORT\_MEASURE\_LEDGER\_1}
\end{Highlighting}
\end{Shaded}

and dominance checker:

\begin{Shaded}
\begin{Highlighting}[]
\NormalTok{YM3D\_DEFECT\_ENTROPY\_DOMINANCE\_CHECK\_1}
\end{Highlighting}
\end{Shaded}

\begin{center}\rule{0.5\linewidth}{0.5pt}\end{center}

\subsection{10. Fallback Basin
Certification}\label{fallback-basin-certification}

The UV ladder does not by itself prove infrared clustering. A separate
row must certify entry into a fallback basin where a strong-coupling,
character-expansion, defect-expansion, transfer-matrix, or terminal
polymer argument supplies volume-uniform decay.

The required artifact is

\begin{Shaded}
\begin{Highlighting}[]
\NormalTok{YM3D\_FALLBACK\_BASIN\_CERTIFICATION\_1}
\end{Highlighting}
\end{Shaded}

It must not be replaced by the weaker handoff row. The handoff row
records that the UV flow has reached the relevant scale; the
certification row proves that the terminal-scale theory lies inside a
convergence region with explicit constants.

A valid certification row must include:

\begin{longtable}[]{@{}
  >{\raggedright\arraybackslash}p{(\columnwidth - 2\tabcolsep) * \real{0.5000}}
  >{\raggedright\arraybackslash}p{(\columnwidth - 2\tabcolsep) * \real{0.5000}}@{}}
\toprule\noalign{}
\begin{minipage}[b]{\linewidth}\raggedright
Field
\end{minipage} & \begin{minipage}[b]{\linewidth}\raggedright
Meaning
\end{minipage} \\
\midrule\noalign{}
\endhead
\bottomrule\noalign{}
\endlastfoot
\texttt{fallback\_method} & character expansion, defect expansion,
polymer expansion, transfer matrix, or hybrid \\
\texttt{terminal\_scale\_id} & scale (k\^{}\ast) entering the fallback
analysis \\
\texttt{input\_activity\_vector} & vector imported from the UV ladder \\
\texttt{convergence\_criterion} & KP, Dobrushin, Fernandez--Procacci, or
other certified criterion \\
\texttt{mass\_scale\_bound} & decay scale or spectral-gap lower bound at
the terminal scale \\
\texttt{volume\_uniformity\_row} & proof that constants are independent
of periodic volume \\
\texttt{overlap\_or\_majorization\_check} & comparison between UV
terminal action and fallback representation \\
\texttt{nonclaims} & no continuum claim unless (a)-uniformity is
separately proved \\
\end{longtable}

The fallback certification gate is

{[}
\mathbf z\_\{3D\}\textsuperscript{\{\rm terminal\}(k}\ast,I\_\{k\^{}\ast\})
\in \mathcal D\_\{\rm fallback\} {]}

with the fallback convergence domain explicitly defined by the chosen
convergence criterion.

\begin{center}\rule{0.5\linewidth}{0.5pt}\end{center}

\subsection{11. Mandatory Sanity Gates}\label{mandatory-sanity-gates}

\subsubsection{Gate S1: cutoff leakage}\label{gate-s1-cutoff-leakage}

For smooth cutoffs, every derivative contribution must be bounded:

{[} C\_\{\rm deriv\}(k,m)p\_k\^{}\{-m\}z\_\{\rm local\}(k,I\_k)
\le R\_\{\rm leakage\}(k,I\_k). {]}

For sharp cutoffs, no derivative may hit a cutoff unless an explicit
boundary-integration row exists. Any such row must include
\texttt{boundary\_surface\_multiplier\_id} and an independently
checkable trace/surface-volume certificate.

\subsubsection{Gate S2: Haar
reversibility}\label{gate-s2-haar-reversibility}

Every transition from gauge-fixed to invariant lane must prove

{[} dU=dU\^{}\{\rm gf\},dh,J\_\{\rm gf\}, \qquad \int dh=1. {]}

\subsubsection{Gate S3: typed compatibility
sparsity}\label{gate-s3-typed-compatibility-sparsity}

If two support types cannot geometrically intersect under the declared
block geometry, then

{[} A\_\{ts\}\^{}\{3D\}=0. {]}

A nonzero entry in such a slot fails the row.

\subsubsection{Gate S4: dimensional-flow monotonicity with gauge-ripple
budget}\label{gate-s4-dimensional-flow-monotonicity-with-gauge-ripple-budget}

Every scale transition must verify

{[} \frac{\beta_k^+}{L}
+\textbar{}\Delta\emph{k\^{}\{\rm analytic\}\textbar{}}\{\max\}
+\textbar{}\Delta\emph{k\^{}\{\rm trunc\}\textbar{}}\{\max\}
+\textbar{}\Delta\emph{k\^{}\{\rm gauge\}\textbar{}}\{\max\}
+\textbar{}\Delta\emph{k\^{}\{\rm tail\}\textbar{}}\{\max\}
\textless{}\beta\_k\^{}-. {]}

The gauge-ripple contribution must be zero by proof or explicitly
bounded. All components in this inequality must be expressed in the same
declared \texttt{flow\_error\_norm}; heterogeneous bounds require
conversion certificates before summation.

\subsubsection{Gate S5: orbit-reduction expanded-graph
audit}\label{gate-s5-orbit-reduction-expanded-graph-audit}

Orbit-reduced counts must agree with fully expanded counts on designated
small audit boxes by explicit projection/isomorphism and multiplicity
checks.

\subsubsection{Gate S5a: orbit-closure
consistency}\label{gate-s5a-orbit-closure-consistency}

For every stored orbit representative and every declared element of
(O\_h), the transformed atom must map back to the same representative
under the canonical representative map. The checker must verify
representative idempotence, group-action closure on the valid atom
universe, and agreement of orbit sizes with stabilizer sizes. This gate
closes before the expanded-graph audit.

\subsubsection{Gate S6: periodic volume
uniformity}\label{gate-s6-periodic-volume-uniformity}

Every thermodynamic claim must prove constants independent of periodic
volume.

\subsubsection{Gate S7: terminal/nonterminal
separation}\label{gate-s7-terminalnonterminal-separation}

Terminal-only rows may not claim full RG smallness.

\subsubsection{Gate S8: exact
arithmetic}\label{gate-s8-exact-arithmetic}

All smallness margins must be stored in exact rational or interval form.
Decimal values are display-only.

\subsubsection{Gate S8b: primitive enclosure
certificates}\label{gate-s8b-primitive-enclosure-certificates}

Any analytic row using a non-rational primitive constant must link to a
\texttt{primitive\_enclosure\_certificate}. This includes
Bessel/character coefficients, lattice Green's functions, heat-kernel
coefficients, spectral constants, quadrature values, and
recurrence-generated constants. The certificate must state the
mathematical origin of the primitive, the rational input interval, the
rational output enclosure, the certification method, and a truncation or
remainder bound that an independent checker can verify using rational
arithmetic.

\subsubsection{Gate S9: vector-valued defect entropy
dominance}\label{gate-s9-vector-valued-defect-entropy-dominance}

The large-field defect-tail ledger must prove

{[}
\left(\mathbf c\_\{\rm tail\}(k,I\_k)-\boldsymbol\mu\emph{\{\rm ent\}\right)\cdot\mathbf m
\ge\delta}\{\rm tail\}\textbar{}\mathbf m\textbar{} {]}

for every admissible support-measure vector (\mathbf m) and every
relevant scale/coupling interval, with the resulting exponential defect
sum included in (R\_\{\rm tail\}(k,I\_k)). The scalar surface-area test
is only a specialization.

\subsubsection{Gate S10: convergence-criterion
declaration}\label{gate-s10-convergence-criterion-declaration}

Every polymer/activity gate must declare its convergence criterion. If
the row uses a tree-graph improvement rather than classical KP, it must
import the relevant tree-graph or Penrose-identity bound and provide
negative controls showing that overbudget vectors are rejected.

\subsubsection{Gate S10b: ghost determinant sign
consistency}\label{gate-s10b-ghost-determinant-sign-consistency}

A determinant-generated or ghost-expanded term may not enter an ordinary
positive polymer activity bound unless a determinant sign or
absolute-majorant certificate proves that this use is admissible. If the
term is handled by a signed/fermionic cluster expansion, the row must
declare the signed convergence criterion separately and must not be
counted as a positive KP activity without transfer.

\subsubsection{Gate S11: fallback-basin
certification}\label{gate-s11-fallback-basin-certification}

The UV ladder may not claim infrared clustering until
\texttt{YM3D\_FALLBACK\_BASIN\_CERTIFICATION\_1} proves that the
terminal-scale effective theory lies in a certified strong-coupling,
character-expansion, defect-expansion, transfer-matrix, or polymer
convergence domain.

\subsubsection{Gate S12: reflection-positivity transfer after gauge
averaging}\label{gate-s12-reflection-positivity-transfer-after-gauge-averaging}

If a row enters an OS-facing claim, the checker must verify the
machine-parseable transition pipeline

\begin{Shaded}
\begin{Highlighting}[]
\NormalTok{gauge\_fixed\_object}
\NormalTok{  {-}\textgreater{} haar\_averaged\_invariant\_object}
\NormalTok{  {-}\textgreater{} reflection\_tagged\_object}
\NormalTok{  {-}\textgreater{} os\_admissible\_object}
\end{Highlighting}
\end{Shaded}

Equivalently, the mathematical transition is

{[} \text{gauge-fixed object} \to \text{Haar-averaged invariant object}
\to \text{reflection-parity-tagged object} \to
\text{OS-admissible observable}. {]}

Gauge-lane metadata and Haar reversibility are necessary but not
sufficient; the final object must also carry an
\texttt{rp\_safe\_use=true} tag imported from a parity/reflection
ledger. In addition, the Haar-averaging support must pass the
reflection-plane factorization test: it must lie strictly in one
half-space relative to the declared reflection plane, or be exactly
reflection-symmetric with a certified pairing of integrated gauge
variables. An asymmetric crossing of the reflection plane fails Gate
S12.

\subsubsection{Gate S13: continuum-uniformity
placeholder}\label{gate-s13-continuum-uniformity-placeholder}

No row may promote volume-uniform lattice clustering to a continuum
mass-gap statement unless an implemented (a)-uniformity ledger proves
that the physical mass lower bound remains positive as (a\to 0). This
gate is a placeholder in the present design, but it is a hard dependency
for any lattice-spacing-uniform physical gap claim.

\begin{center}\rule{0.5\linewidth}{0.5pt}\end{center}

\subsection{12. Operational Milestones}\label{operational-milestones}

\subsubsection{Milestone 0: Exact two-dimensional SU(N) validation
lane}\label{milestone-0-exact-two-dimensional-sun-validation-lane}

Milestone 0 is achieved when
\texttt{YM2D\_SUN\_EXACT\_VALIDATION\_LANE\_0} runs the ledger pipeline
on a two-dimensional SU(N) benchmark and compares its observable output
against exact Wilson-loop/area-law formulas.

Milestone 0 is not required before the three-dimensional combinatorial
Milestone 1A generator. It is required before analytic three-dimensional
rows close.

Allowed claim:

{[}
\text{The ledger pipeline reproduces the declared exact 2D benchmark observables.}
{]}

Not allowed:

{[} \text{3D analytic smallness, RG contraction, or mass-gap claims.}
{]}

\subsubsection{Milestone 1: Pure combinatorial
closure}\label{milestone-1-pure-combinatorial-closure}

Closed rows:

\begin{itemize}
\tightlist
\item
  \texttt{YM3D\_LOCAL\_ATOM\_GEOMETRY\_1}
\item
  \texttt{YM3D\_BLOCK\_GAUGE\_FIXING\_LEDGER\_1}
\item
  \texttt{YM3D\_BALABAN\_TREE\_AXIAL\_GAUGE\_HOOK\_1}
\item
  \texttt{YM3D\_HAAR\_REVERSIBILITY\_GAUGE\_CHECK\_1}
\item
  \texttt{YM3D\_CUBIC\_SYMMETRY\_ORBIT\_REDUCTION\_1}
\item
  \texttt{YM3D\_ORBIT\_REDUCTION\_EXPANDED\_GRAPH\_AUDIT\_1}
\end{itemize}

Meaning: the local geometry, gauge trees, residual gauge metadata, orbit
representatives, support types, stabilizer multipliers, and
reduced-vs-expanded collision graph multiplicities are verified. No
analytic smallness has been proved.

\subsubsection{Milestone 1A: Single-block executable
audit}\label{milestone-1a-single-block-executable-audit}

This is the first executable checkpoint after pure design. On one or
more small periodic audit boxes, the generator must produce both:

\begin{enumerate}
\def\labelenumi{\arabic{enumi}.}
\tightlist
\item
  the fully expanded atom/collision graph; and
\item
  the orbit-reduced graph with stabilizer and multiplicity data.
\end{enumerate}

The independent checker must verify the exact projection identity

{[} M\^{}\{\rm expanded\}\emph{\{{[}\alpha{]},{[}\beta{]},q\} =
M\^{}\{\rm reduced\}}\{{[}\alpha{]},{[}\beta{]},q\} {]}

for every orbit-pair collision type (q).

Milestone 1A is the first target for automated testing. It is
deliberately decoupled from analytic smallness.

\subsubsection{Milestone 2: Single-scale UV
bound}\label{milestone-2-single-scale-uv-bound}

For the first weak-coupling slab (I\_0), the full scale-indexed activity
gate passes:

{[} \sum\emph{s A}\{ts\}\^{}\{3D\}(0)
\left(z\_s\textsuperscript{\{\rm term\}(0,I\_0)+z\_s}\{\rm nonterm\}(0,I\_0)\right)e\^{}\{a\_s(0)\}
\le a\_t(0). {]}

Meaning: one UV scale has local KP stability.

\subsubsection{Milestone 3: Multi-scale UV
ladder}\label{milestone-3-multi-scale-uv-ladder}

A sequence

{[} I\_0\to I\_1\to\cdots\to I\_K {]}

passes both dimensional-flow monotonicity and the full activity gate at
each scale.

\subsubsection{Milestone 4: Fallback basin
certification}\label{milestone-4-fallback-basin-certification}

The ladder reaches a threshold (k\^{}\ast) where

{[} g\_3\^{}2 a\_\{k\^{}\ast\}\asymp 1 {]}

or (I\_\{k\^{}\ast\}) lies in a certified fallback interval. The proof
then transitions from the UV Taylor/small-field map to a
strong-coupling, character-expansion, defect-expansion, transfer-matrix,
or terminal polymer map, and
\texttt{YM3D\_FALLBACK\_BASIN\_CERTIFICATION\_1} proves that the
terminal-scale effective theory lies inside the declared convergence
domain.

\subsubsection{Milestone 5: Volume-uniform
clustering}\label{milestone-5-volume-uniform-clustering}

The cluster/polymer expansion proves

{[} \textbar{}\langle O(x)O(y)\rangle\_c\textbar{}
\le C\_O(k)e\^{}\{-m\_k d(x,y)\} {]}

with constants independent of finite periodic volume.

\subsubsection{Milestone 6: Lattice-spacing-uniform physical gap
window}\label{milestone-6-lattice-spacing-uniform-physical-gap-window}

This milestone is closed only after Gate S13 is implemented and passed.
The required result is that the clustering scale yields

{[} m\_\{\rm phys\}\ge c g\_3\^{}2, \qquad c\textgreater0, {]}

uniformly in the continuum scaling window. Before Gate S13 closes, the
strongest allowed claim is volume-uniform lattice clustering inside
certified slabs, not a continuum-stable physical mass lower bound.
Milestone 6 uses the default sequential limit order: first
(\Lambda  o\infty) at fixed (a), then (a o0) under Gate S13 uniformity.

\begin{center}\rule{0.5\linewidth}{0.5pt}\end{center}

\subsection{13. Artifact Sequence}\label{artifact-sequence}

The recommended first sequence is:

\begin{enumerate}
\def\labelenumi{\arabic{enumi}.}
\setcounter{enumi}{-1}
\tightlist
\item
  \texttt{YM3D\_SCHEMA\_STAGE\_DISCRIMINATOR\_1} 0a.
  \texttt{YM3D\_SCHEMA\_DISCRIMINATOR\_NULL\_TEST\_1} 0b.
  \texttt{YM3D\_PIPELINE\_DAG\_DECLARATION\_1}
\item
  \texttt{YM3D\_LOCAL\_ATOM\_GEOMETRY\_1}
\item
  \texttt{YM3D\_BLOCK\_GAUGE\_FIXING\_LEDGER\_1}
\item
  \texttt{YM3D\_BALABAN\_TREE\_AXIAL\_GAUGE\_HOOK\_1}
\item
  \texttt{YM3D\_HAAR\_REVERSIBILITY\_GAUGE\_CHECK\_1}
\item
  \texttt{YM3D\_CUBIC\_SYMMETRY\_ORBIT\_REDUCTION\_1} 5a.
  \texttt{YM3D\_ORBIT\_CLOSURE\_CONSISTENCY\_CHECK\_1}
\item
  \texttt{YM3D\_ORBIT\_REDUCTION\_EXPANDED\_GRAPH\_AUDIT\_1}
\item
  \texttt{YM3D\_SCALE\_COUPLING\_SLAB\_LEDGER\_1} 7a.
  \texttt{YM3D\_FLOW\_ERROR\_NORM\_LEDGER\_1}
\item
  \texttt{YM3D\_SMALL\_LARGE\_FIELD\_PARTITION\_LEDGER\_1}
\item
  \texttt{YM3D\_CUTOFF\_LEAKAGE\_CHECK\_1}
\item
  \texttt{YM3D\_KSTAR\_TERMINAL\_ATOM\_EXPORT\_1}
\item
  \texttt{YM3D\_KSTAR\_TERMINAL\_COLLISION\_GRAPH\_1}
\item
  \texttt{YM3D\_KP\_SPARSITY\_GEOMETRY\_CHECK\_1}
\item
  \texttt{YM3D\_TERMINAL\_KP\_ACTIVITY\_MAP\_SCALE\_INDEXED\_1}
\item
  \texttt{YM3D\_NONTERMINAL\_TAYLOR\_TERM\_INVENTORY\_1}
\item
  \texttt{YM3D\_FP\_DETERMINANT\_EXPANSION\_LEDGER\_1} 15a.
  \texttt{YM3D\_GHOST\_DETERMINANT\_SIGN\_CONSISTENCY\_CHECK\_1}
\item
  \texttt{YM3D\_BOUNDARY\_SURFACE\_MULTIPLIER\_LEDGER\_1}
\item
  \texttt{YM3D\_PRIMITIVE\_ENCLOSURE\_CERTIFICATE\_LEDGER\_1}
\item
  \texttt{YM3D\_RP\_SPATIAL\_FACTORIZATION\_CHECK\_1}
\item
  \texttt{YM3D\_GAUGE\_RIPPLE\_FLOW\_CORRECTION\_1}
\item
  \texttt{YM3D\_THERMODYNAMIC\_DEFECT\_TAIL\_LEDGER\_1}
\item
  \texttt{YM3D\_DEFECT\_SUPPORT\_MEASURE\_LEDGER\_1}
\item
  \texttt{YM3D\_DEFECT\_ENTROPY\_DOMINANCE\_CHECK\_1}
\item
  \texttt{YM3D\_BKAR\_FOREST\_REMAINDER\_LEDGER\_1} if forest
  interpolation is used
\item
  \texttt{YM3D\_ONE\_STEP\_RG\_RUNG\_SCALE\_INDEXED\_1}
\item
  \texttt{YM3D\_FALLBACK\_BASIN\_HANDOFF\_1}
\item
  \texttt{YM3D\_FALLBACK\_BASIN\_CERTIFICATION\_1}
\item
  \texttt{YM3D\_VOLUME\_UNIFORM\_CLUSTERING\_1}
\item
  \texttt{YM3D\_A\_UNIFORM\_GAP\_WINDOW\_1} Validation and advisory
  lanes are maintained outside the critical three-dimensional proof DAG
  unless a formal transfer row is added:
\end{enumerate}

\begin{itemize}
\tightlist
\item
  \texttt{YM2D\_SUN\_EXACT\_VALIDATION\_LANE\_0} -- mandatory before
  analytic three-dimensional rows close;
\item
  \texttt{YM3D\_NUMERICAL\_ANCHOR\_LANE\_0} -- advisory
  numerical/source-audit lane;
\item
  \texttt{YM3D\_OBSERVABLE\_REGISTRY\_1} -- mandatory registry for
  observable-specific claims;
\item
  \texttt{YM3D\_PHI4\_BENCHMARK\_LANE\_0} -- advisory scalar benchmark;
\item
  \texttt{YM3D\_U1\_BENCHMARK\_LANE\_0} -- advisory compact abelian
  benchmark;
\item
  \texttt{YM3D\_SPDE\_COMPARISON\_LANE\_0} -- advisory
  stochastic-quantization comparison lane;
\item
  \texttt{YM3D\_HAMILTONIAN\_COMPARISON\_LANE\_0} -- advisory
  Hamiltonian/orbit-space comparison lane;
\item
  \texttt{YM3D\_FORMAL\_FINITE\_GROUP\_AUDIT\_LANE\_1} -- formalization
  lane for stable finite group-action/orbit invariants after Milestone
  1A.
\end{itemize}

\begin{center}\rule{0.5\linewidth}{0.5pt}\end{center}

\subsection{14. Checker Requirements}\label{checker-requirements}

Every checker must include positive and negative controls. The design
distinguishes hard mathematical negative controls from routine
schema-hygiene checks. The document foregrounds hard controls;
duplicated identifiers, missing parent IDs, malformed references, and
similar hygiene failures belong in CI/checker documentation.

Hard negative controls include:

\begin{itemize}
\tightlist
\item
  illegal signed cancellation;
\item
  finite-volume claim promoted to thermodynamic or continuum claim;
\item
  terminal-only row promoted to full-RG row;
\item
  gauge-fixed row used in invariant claim without Haar transfer;
\item
  asymmetric Haar/RP support crossing a reflection plane;
\item
  smooth cutoff derivative leakage omitted;
\item
  sharp cutoff differentiated without a boundary-surface multiplier
  ledger;
\item
  nonzero compatibility entry where geometry requires zero;
\item
  nonmonotone (\beta\_k) flow accepted;
\item
  heterogeneous (\Delta\_k) norms summed without conversion
  certificates;
\item
  gauge-ripple correction omitted from the flow-error budget;
\item
  vector-valued defect entropy rate not dominated by the tail
  action-cost vector;
\item
  scalar surface-area defect model used for a lower-dimensional,
  junctioned, or localized high-curvature defect;
\item
  missing junction/branching entropy component for a mixed defect
  complex;
\item
  orbit-reduction multipliers not verified against expanded audit
  graphs;
\item
  orbit representative fails closure under an element of \texttt{O\_h};
\item
  representative map is not idempotent;
\item
  composite support orbit action does not match the lift from
  transformed constituent links;
\item
  bounded-field approximation reported as full large-field control;
\item
  fallback-basin handoff reported as fallback-basin certification;
\item
  terminal scalar bound used without typed vector export;
\item
  polymer criterion changed without updating the convergence-domain
  certificate;
\item
  ghost-dependent determinant term missing a ghost-measure certificate;
\item
  ghost-expanded determinant term entering a positive polymer bound
  without sign/absolute-majorant admissibility;
\item
  configuration-dependent determinant compressed into an untyped static
  coefficient;
\item
  primitive enclosure interval accepted without an independently
  checkable remainder certificate;
\item
  primitive enclosure backend accepted without version, precision,
  directed-rounding, or rationalized output metadata;
\item
  comparison-lane or numerical-anchor result imported as a proof
  dependency without a formal transfer row;
\item
  tree/block-axial gauge lane propagated into fallback or clustering
  claims without gauge-covariant or invariant transfer;
\item
  Milestone 6 claimed before Gate S13 is implemented and closed;
\item
  analytic fields accepted as nullable in a combinatorial schema branch
  rather than forbidden by a stage discriminator;
\item
  pipeline DAG dependency missing or inconsistent with row parent IDs;
\item
  discriminator null-test fixtures do not produce the expected pass/fail
  pattern;
\item
  observable-specific clustering claim made without an
  observable-registry entry;
\item
  claim-strength upgrade made without passing the corresponding DAG node
  and checker;
\item
  analytic three-dimensional row closed before the exact two-dimensional
  validation lane has passed or an explicit waiver has been recorded.
\end{itemize}

\begin{center}\rule{0.5\linewidth}{0.5pt}\end{center}

\subsection{15. Automation and Formalization
Strategy}\label{automation-and-formalization-strategy}

The first implementation should prioritize executable correctness over
full proof-assistant formalization. The first executable target is not
atom generation; it is the discriminator null-test harness:

\begin{Shaded}
\begin{Highlighting}[]
\NormalTok{validate\_schema\_discriminators.py}
\end{Highlighting}
\end{Shaded}

Only after this harness passes should
\texttt{YM3D\_LOCAL\_ATOM\_GEOMETRY\_1} emit concrete atoms.

The staged stack is:

\begin{enumerate}
\def\labelenumi{\arabic{enumi}.}
\tightlist
\item
  Python exact-arithmetic generator for finite geometry, group actions,
  orbit representatives, and expanded audit graphs.
\item
  JSON Schema or equivalent machine-readable schemas for all row types.
\item
  Independent checker scripts that do not share internal generator code.
\item
  Property-based negative controls for orbit counts, gauge-lane misuse,
  illegal support intersections, and missing stabilizers.
\item
  Lean 4 or another proof assistant for the finite
  group-action/orbit/stabilizer invariants immediately after Milestone
  1A stabilizes the object vocabulary. This formal lane should cover
  \texttt{O\_h} action closure, representative idempotence,
  stabilizer/orbit identities, and canonical lift consistency.
\end{enumerate}

The repository should treat generated ledgers as artifacts and checkers
as independent validators. A generator success is not a proof; only an
independent checker pass closes a row.

The schema layer must use discriminated unions keyed by
\texttt{schema\_stage} or \texttt{milestone\_target}; milestone-specific
enforcement is not deferred to Python alone. The Python checker verifies
semantic mathematics, but the JSON Schema must reject structurally
invalid rows immediately.

Recommended schema families:

\begin{longtable}[]{@{}
  >{\raggedright\arraybackslash}p{(\columnwidth - 2\tabcolsep) * \real{0.5000}}
  >{\raggedright\arraybackslash}p{(\columnwidth - 2\tabcolsep) * \real{0.5000}}@{}}
\toprule\noalign{}
\begin{minipage}[b]{\linewidth}\raggedright
Schema
\end{minipage} & \begin{minipage}[b]{\linewidth}\raggedright
Purpose
\end{minipage} \\
\midrule\noalign{}
\endhead
\bottomrule\noalign{}
\endlastfoot
\texttt{stage\_discriminated\_row.schema.json} & common \texttt{oneOf}
discriminator pattern for combinatorial, gauge-geometric, analytic,
RP-transfer, and continuum-uniformity rows \\
\texttt{pipeline\_dag.schema.json} & machine-readable proof/data-flow
DAG with node dependencies and claims unlocked \\
\texttt{convergence\_criterion.schema.json} & KP, Dobrushin, tree-graph,
signed/fermionic, and determinant-admissibility convergence modes \\
\texttt{atom.schema.json} & local atom/support/gauge metadata \\
\texttt{orbit.schema.json} & cubic orbit representatives, stabilizers,
and upward composition rules from links to
plaquettes/cubes/multiblocks \\
\texttt{collision.schema.json} & expanded and reduced collision edges \\
\texttt{nonterminal\_term.schema.json} & Taylor, forest, determinant,
ghost-measure, cutoff, and primitive-certificate fields for analytic
terms \\
\texttt{activity\_vector.schema.json} & typed terminal/nonterminal
activity vectors \\
\texttt{convergence\_criterion.schema.json} & KP/Dobrushin/tree-graph
domain certificates \\
\texttt{defect\_support.schema.json} & vector-valued support measures,
junction/branching data, and entropy dominance data \\
\texttt{defect\_entropy\_certificate.schema.json} & action-cost/entropy
dominance certificates for mixed defects \\
\texttt{fp\_determinant\_certificate.schema.json} &
configuration-dependent FP or determinant-expansion certificates \\
\texttt{rp\_transfer.schema.json} & reflection-plane factorization and
OS-transfer certificates \\
\texttt{boundary\_surface\_certificate.schema.json} & sharp-cutoff
boundary surface/trace multiplier certificates \\
\texttt{primitive\_enclosure\_certificate.schema.json} & rationally
checkable enclosures for non-rational analytic constants \\
\texttt{scale\_coupling\_slab.schema.json} & scale slabs, common
flow-error norm, converted (\Delta\_k) components, and monotonicity
claims \\
\texttt{continuum\_limit\_policy.schema.json} & thermodynamic/continuum
limit-order policy and Gate S13 dependencies \\
\texttt{milestone.schema.json} & closed rows and non-claims per
milestone \\
\end{longtable}

Base geometry rows do not carry nullable analytic certificate fields.
Instead, their discriminated schema branch forbids analytic-only fields.
Later analytic branches require the relevant certificate IDs. This
preserves Milestone 1A as a combinatorial task while preventing both
breaking schema changes and accidental acceptance of incomplete analytic
rows.

Rapid prototyping may use a separate
\texttt{schema\_profile:\ "experimental"} branch, but experimental
artifacts cannot close gates. Gate-closing artifacts must use
\texttt{schema\_profile:\ "strict"} and the stage-discriminated
production schemas.

\begin{center}\rule{0.5\linewidth}{0.5pt}\end{center}

\subsection{16. Validation, Benchmark, Numerical, and Comparison
Lanes}\label{validation-benchmark-numerical-and-comparison-lanes}

The main proof DAG remains the lattice constructive RG path. Validation
and comparison lanes do not close proof gates unless a formal transfer
row is added. They are nevertheless required to prevent the
schema/checker machinery from drifting away from known mathematics and
external data.

\subsubsection{16.1 Exact two-dimensional SU(N) validation
lane}\label{exact-two-dimensional-sun-validation-lane}

\texttt{YM2D\_SUN\_EXACT\_VALIDATION\_LANE\_0} is the strongest
validation lane. It runs the combinatorial, orbit-reduction, observable,
and Wilson-loop portions of the ledger pipeline on two-dimensional SU(N)
Yang--Mills and compares output against exact Wilson-loop/area-law
formulas.

This lane is not a prerequisite for Milestone 1A, but it blocks analytic
three-dimensional rows unless an explicit waiver is recorded. The lane
is intended to catch errors in link-to-plaquette lifting, orbit
multipliers, gauge-lane transfers, reflection-factorization metadata,
and observable normalization before the three-dimensional
Taylor/defect/BKAR rows are attempted.

\subsubsection{16.2 Numerical anchor lane}\label{numerical-anchor-lane}

\texttt{YM3D\_NUMERICAL\_ANCHOR\_LANE\_0} compares generator outputs
against small-volume exact enumeration where possible and against
external Monte Carlo or high-precision lattice data where appropriate.
Numerical anchors may compare:

\begin{itemize}
\tightlist
\item
  terminal activity vector components;
\item
  small-field Taylor coefficients;
\item
  Wilson-loop estimates;
\item
  scalar glueball correlator proxies;
\item
  string-tension or glueball-mass ratios used only as external sanity
  checks.
\end{itemize}

Numerical anchors never prove a theorem. They emit source-audit warnings
when analytic normalizations, coefficient bounds, or gauge-ripple
estimates disagree with independent numerical evidence.

\subsubsection{16.3 Observable registry}\label{observable-registry}

\texttt{YM3D\_OBSERVABLE\_REGISTRY\_1} is mandatory for any claim beyond
pure combinatorics. Each analytic, clustering, fallback, or
continuum-uniformity row must declare which registered gauge-invariant
observables it controls.

The initial registry contains Wilson loops, scalar glueball operators,
connected scalar glueball correlators, and Polyakov-loop diagnostics.
Polyakov loops are finite-temperature diagnostics and do not by
themselves constitute a zero-temperature mass-gap claim.

\subsubsection{16.4 Scalar benchmark lane}\label{scalar-benchmark-lane}

\texttt{YM3D\_PHI4\_BENCHMARK\_LANE\_0} tests polymer/KP, cutoff
leakage, forest formulas, and periodic lifts without gauge-lane
complexity. It should not delay Milestone 1 for Yang--Mills.

\subsubsection{16.5 Abelian compact gauge benchmark
lane}\label{abelian-compact-gauge-benchmark-lane}

\texttt{YM3D\_U1\_BENCHMARK\_LANE\_0} tests compact gauge variables,
defect tails, and entropy suppression against known three-dimensional
compact abelian confinement structures. It is a benchmark for tail
accounting, not a substitute for non-Abelian control.

\subsubsection{16.6 Stochastic quantization comparison
lane}\label{stochastic-quantization-comparison-lane}

\texttt{YM3D\_SPDE\_COMPARISON\_LANE\_0} tracks continuum
stochastic-quantization work as a comparison route. It may inform
lattice-to-continuum observables, gauge-covariant state spaces,
renormalization metadata, or quotient-space design, but it is not a
dependency of the lattice proof DAG.

\subsubsection{16.7 Hamiltonian/orbit-space comparison
lane}\label{hamiltonianorbit-space-comparison-lane}

\texttt{YM3D\_HAMILTONIAN\_COMPARISON\_LANE\_0} tracks
Karabali--Kim--Nair-type Hamiltonian/orbit-space approaches in (2+1)
dimensions. This lane may suggest relevant operators or mass scales, but
it does not replace Euclidean lattice RG certification.

\subsubsection{16.8 Formal finite-group audit
lane}\label{formal-finite-group-audit-lane}

\texttt{YM3D\_FORMAL\_FINITE\_GROUP\_AUDIT\_LANE\_1} begins after
Milestone 1A stabilizes the finite object vocabulary. It should
formalize the cubic group action, orbit-closure theorem,
stabilizer/orbit identity, canonical representative idempotence, and
upward composition of link actions to composite supports. The initial
production generator remains Python exact-arithmetic; the formal lane
prevents finite-combinatorial drift after the objects are known.

\begin{center}\rule{0.5\linewidth}{0.5pt}\end{center}

\subsection{17. Literature Anchors}\label{literature-anchors}

The design uses the following established methods as anchors or
cautionary templates.

\subsubsection{Balaban ultraviolet
stability}\label{balaban-ultraviolet-stability}

Balaban proved ultraviolet stability for three-dimensional Wilson
lattice pure Yang--Mills theories using a renormalization-group method
for lattice gauge theories. This is the closest methodological anchor
for the present program.

\subsubsection{Tree-axial gauge methods}\label{tree-axial-gauge-methods}

Balaban-style block gauge fixing motivates the explicit tree-axial gauge
ledger. The program must treat gauge fixing, residual gauge variables,
and Haar reversibility as first-class proof objects.

\subsubsection{Göpfert--Mack three-dimensional U(1)
confinement}\label{guxf6pfertmack-three-dimensional-u1-confinement}

Göpfert and Mack proved confinement of static quarks in
three-dimensional compact (U(1)) lattice gauge theory for all couplings.
This is not a direct non-Abelian (SU(2)) substitute. It is used here as
a structural guide for thermodynamic defect-tail accounting.

\subsubsection{Dimock/Balaban RG expositions and bounded-field
approximations}\label{dimockbalaban-rg-expositions-and-bounded-field-approximations}

Dimock's expositions and applications of Balaban-style RG methods
clarify the role of bounded-field approximations and the danger of
confusing a bounded-field result with full large-field control.

\subsubsection{BKAR/forest formulas}\label{bkarforest-formulas}

If nonterminal remainders are organized by forest interpolation, the
Brydges--Kennedy--Abdesselam--Rivasseau family of forest formulas must
be declared explicitly. The formula determines multiblock combinatorics
and compatibility types.

\subsubsection{Abstract polymer convergence
improvements}\label{abstract-polymer-convergence-improvements}

Classical KP is the default starting point, but improved abstract
polymer criteria, including Dobrushin-type and
Fernandez--Procacci/tree-graph bounds, should be available as
alternative convergence domains when absolute KP majorization is too
crude.

\subsubsection{Stochastic quantization and SPDE
comparison}\label{stochastic-quantization-and-spde-comparison}

Recent stochastic-quantization constructions for three-dimensional
Yang--Mills--Higgs theories provide a parallel analytic context
involving gauge-covariant stochastic flows and Markov processes on gauge
orbits. These results are comparison material, not proof dependencies.

\subsubsection{Hamiltonian/orbit-space
approaches}\label{hamiltonianorbit-space-approaches}

Karabali--Kim--Nair-type Hamiltonian methods in (2+1)-dimensional
Yang--Mills supply independent physical and analytic intuition for mass
generation. They may guide operator selection but are not used as a
Euclidean lattice RG proof substitute.

\subsubsection{Two-dimensional Yang--Mills exact
validation}\label{two-dimensional-yangmills-exact-validation}

Two-dimensional Yang--Mills provides an exact validation target for the
ledger machinery. The discrete and continuum constructions on compact
surfaces provide exact Wilson-loop and area-law comparison data against
which orbit lifting, gauge normalization, and observable export can be
tested before three-dimensional analytic rows close.

\subsubsection{Numerical lattice anchors in 2+1
dimensions}\label{numerical-lattice-anchors-in-21-dimensions}

Lattice calculations of SU(N) gauge theories in (2+1) dimensions provide
external numerical anchors for glueball spectra, string tensions, and
dimensionless continuum ratios. These data are not proof inputs, but
they are useful source-audit checks for normalization and
coefficient-bound sanity.

\begin{center}\rule{0.5\linewidth}{0.5pt}\end{center}

\subsection{18. Minimal Repository
Layout}\label{minimal-repository-layout}

\begin{Shaded}
\begin{Highlighting}[]
\NormalTok{ym3d{-}constructive{-}rg/}
\NormalTok{  README.md}
\NormalTok{  YM3D\_CONSTRUCTIVE\_RG\_DESIGN.md}
\NormalTok{  references.bib}
\NormalTok{  ledgers/}
\NormalTok{    README.md}
\NormalTok{  schemas/}
\NormalTok{    README.md}
\NormalTok{    stage\_discriminated\_row.schema.json}
\NormalTok{    atom.schema.json}
\NormalTok{    orbit.schema.json}
\NormalTok{    collision.schema.json}
\NormalTok{    nonterminal\_term.schema.json}
\NormalTok{    activity\_vector.schema.json}
\NormalTok{    convergence\_criterion.schema.json}
\NormalTok{    observable\_registry.schema.json}
\NormalTok{    numerical\_anchor.schema.json}
\NormalTok{    ym2d\_exact\_validation.schema.json}
\NormalTok{    defect\_support.schema.json}
\NormalTok{    defect\_entropy\_certificate.schema.json}
\NormalTok{    fp\_determinant\_certificate.schema.json}
\NormalTok{    rp\_transfer.schema.json}
\NormalTok{    boundary\_surface\_certificate.schema.json}
\NormalTok{    primitive\_enclosure\_certificate.schema.json}
\NormalTok{    scale\_coupling\_slab.schema.json}
\NormalTok{    continuum\_limit\_policy.schema.json}
\NormalTok{  checkers/}
\NormalTok{    README.md}
\NormalTok{  generators/}
\NormalTok{    README.md}
\NormalTok{  artifacts/}
\NormalTok{    README.md}
\NormalTok{  lanes/}
\NormalTok{    ym2d\_sun\_exact\_validation/}
\NormalTok{    numerical\_anchor/}
\NormalTok{    observable\_registry/}
\NormalTok{    phi4\_benchmark/}
\NormalTok{    u1\_benchmark/}
\NormalTok{    spde\_comparison/}
\NormalTok{    hamiltonian\_comparison/}
\NormalTok{    formal\_finite\_group\_audit/}
\end{Highlighting}
\end{Shaded}

This initial bundle provides the top-level design and bibliography only.
The subdirectories are reserved for future machine-readable ledgers,
schemas, checkers, and generated artifacts.

\begin{center}\rule{0.5\linewidth}{0.5pt}\end{center}

\subsection{19. Schema-Freeze and Validation Additions in Versions
1.5--1.7}\label{schema-freeze-and-validation-additions-in-versions-1.51.7}

Version 1.5 adds the following schema-freeze requirements before
Milestone 1A implementation:

\begin{enumerate}
\def\labelenumi{\arabic{enumi}.}
\tightlist
\item
  discriminated schema branches keyed by \texttt{schema\_stage} or
  \texttt{milestone\_target};
\item
  forbidden analytic fields in purely combinatorial rows rather than
  nullable analytic placeholders;
\item
  upward composition rules for the (O\_h) action from oriented links to
  plaquettes, cubes, and multiblocks;
\item
  determinant/ghost kinematics, including \texttt{ghost\_measure\_id},
  \texttt{ghost\_cutoff\_policy}, and optional
  \texttt{grassmann\_cutoff\_id};
\item
  norm-typed flow-error budgets, with all
  (\textbar{}\Delta\emph{k\textbar{}}\{\max\}) components expressed in a
  common \texttt{flow\_error\_norm};
\item
  conversion certificates for heterogeneous norms;
\item
  explicit thermodynamic-before-continuum limit order for Gate S13;
\item
  a \texttt{continuum\_limit\_policy.schema.json} recording whether a
  row uses sequential or joint uniformity.
\end{enumerate}

The rule for schema freeze is now:

{[} \text{JSON Schema rejects structurally incomplete rows;} \qquad
\text{independent checkers reject mathematically invalid rows.} {]}

The previous v1.4 schema hooks remain active: mixed-defect topology
fields, determinant metadata, reflection-transfer geometry,
sharp-boundary multiplier certificates, and primitive enclosure
certificates.

\begin{center}\rule{0.5\linewidth}{0.5pt}\end{center}

Version 1.6 adds four implementation-readiness controls:

\begin{enumerate}
\def\labelenumi{\arabic{enumi}.}
\tightlist
\item
  \texttt{YM3D\_ORBIT\_CLOSURE\_CONSISTENCY\_CHECK\_1}, requiring
  closure of representative maps under every declared element of
  \texttt{O\_h} before expanded-graph audits.
\item
  \texttt{YM3D\_GHOST\_DETERMINANT\_SIGN\_CONSISTENCY\_CHECK\_1},
  preventing ghost-expanded determinant terms from entering positive
  activity criteria without signed-cluster or absolute-majorant
  certification.
\item
  \texttt{docs/pipeline\_dag.mmd} and
  \texttt{artifacts/pipeline\_dag.json}, making the proof/data-flow
  sequence human-readable and machine-auditable.
\item
  \texttt{YM3D\_SCHEMA\_DISCRIMINATOR\_NULL\_TEST\_1}, requiring
  pass/fail fixtures for the stage-discriminated schema branches before
  generator output.
\end{enumerate}

Version 1.7 adds validation and claim-anchoring controls:

\begin{enumerate}
\def\labelenumi{\arabic{enumi}.}
\tightlist
\item
  \texttt{YM2D\_SUN\_EXACT\_VALIDATION\_LANE\_0}, mandatory before
  analytic three-dimensional rows close unless explicitly waived;
\item
  \texttt{YM3D\_NUMERICAL\_ANCHOR\_LANE\_0}, an advisory source-audit
  lane comparing generator output to small-box or lattice numerical
  data;
\item
  \texttt{YM3D\_OBSERVABLE\_REGISTRY\_1}, requiring observable-specific
  claim targets for clustering, fallback, and continuum-uniformity rows;
\item
  \texttt{claim\_strength} annotations in \texttt{pipeline\_dag.json},
  distinguishing combinatorial, volume-uniform, \texttt{a}-uniform,
  continuum-subsequential, continuum-mass-gap-candidate, and advisory
  claims;
\item
  explicit DAG nodes for OS reconstruction status and the
  thermodynamic-before-continuum limit-order policy;
\item
  extended primitive-enclosure backend metadata for pure rational, Arb,
  MPFI, Lean-verified, and custom interval enclosures;
\item
  a formal finite-group audit lane for stable \texttt{O\_h}
  orbit/stabilizer/composition invariants after Milestone 1A.
\end{enumerate}

\subsection{20. References}\label{references}

See \texttt{references.bib} for machine-readable BibTeX entries. Core
references include Balaban's 1985 ultraviolet-stability paper,
Göpfert--Mack's 1982 three-dimensional (U(1)) confinement paper,
Dimock's expository Balaban-RG papers and 3D bounded-field RG
applications, Abdesselam--Rivasseau's forest-formula treatment of
cluster expansions, Fernandez--Procacci tree-graph improvements for
abstract polymer models, stochastic-quantization work of
Chandra--Chevyrev--Hairer--Shen, and Karabali--Kim--Nair
Hamiltonian/orbit-space work in (2+1)-dimensional Yang--Mills theory,
Thierry Levy's two-dimensional Yang--Mills construction, and
Athenodorou--Teper/Teper numerical lattice data for SU(N) gauge theories
in (2+1) dimensions.

\end{document}
